\documentclass[10pt,a4paper]{amsart}
\usepackage[margin=19mm]{geometry}
\usepackage{amsmath,amssymb,mathtools}
\usepackage[scr=rsfs,cal=euler]{mathalfa}
\usepackage{xfrac}
\usepackage[unicode,hidelinks,bookmarks=false]{hyperref}
\newcommand{\ie}{i.e.\ }
\newcommand{\cf}{cf.\ }
\newcommand{\resp}{respectively\ }
\newcommand{\Subsection}{\subsection}
\providecommand{\nobreakdash}{\nobreak\hbox{-}}
\setlength{\emergencystretch}{2em}
\multlinegap0pt
\usepackage{lmodern}
\usepackage{mathtools}
\usepackage{booktabs,multirow,array}
\usepackage{tabularx}
\usepackage{float}
\usepackage{placeins}
\usepackage{enumitem}
\allowdisplaybreaks
\newtheorem{theorem}{Theorem}[section]
\newtheorem{proposition}[theorem]{Proposition}
\newtheorem{lemma}[theorem]{Lemma}
\newtheorem{corollary}[theorem]{Corollary}
\newtheorem{definition}[theorem]{Definition}
\newtheorem{remark}[theorem]{Remark}
\newcommand{\R}{\mathbb{R}}
\newcommand{\E}{\mathbb{E}}
\newcommand{\cO}{\mathcal{O}}
\newcommand{\cG}{\mathcal{G}}
\newcommand{\cS}{\mathcal{S}}
\newcommand{\Tr}{\operatorname{Tr}}
\newcommand{\osc}{\operatorname{osc}}
\newcommand{\dist}{\operatorname{dist}}
\newcommand{\dt}{\Delta t}
\begin{document}
\title[A Positivity-Preserving Expectation Scheme for Hamilton--Jac]{A Positivity-Preserving Expectation Scheme for Hamilton--Jacobi--Bellman Equations with Oblique Robin Boundary Conditions}
\author{Haoran Xu$^1$; Haoran Xu, Xingye Yue}
\date{}
\maketitle
\noindent\emph{Source and licence.} A Positivity-Preserving Expectation Scheme for Hamilton--Jacobi--Bellman Equations with Oblique Robin Boundary Conditions. Version arXiv:2608.11936v1 (CC BY 4.0). This material is distributed under the Creative Commons Attribution 4.0 International licence, http://creativecommons.org/licenses/by/4.0/, as recorded in the arXiv OAI licence metadata for this exact version and shown on the abstract page https://arxiv.org/abs/2608.11936v1. Full rights notice: licenses/arxiv-2608.11936-CC-BY-4.0.txt.\par\smallskip
\noindent\textit{Editorial scope.} Counted unit: the original \texttt{theorem} environment \texttt{thm:rate} at source lines 2050--2105 together with its complete proof at lines 2107--2297; the delivered document keeps source lines 238--2417 so that every referenced label and every citation resolves inside it. Native editable source is the paper's own concatenated TeX; the AMS wrapper is editorial formatting only. No publisher class, upstream script or unrelated artwork is redistributed.
	\begin{equation}
		\begin{cases}
			-\partial_t u + H(t,x,Du,D^2u,u) = 0, & (t,x)\in\cO_T, \\[8pt]
			L(t,x,Du,u) = 0, & (t,x)\in[0,T)\times\partial\cO, \\[8pt]
			u(T,x) = \Psi(x), & x\in\overline{\cO},
		\end{cases}
		\label{eq:HJB-main}
	\end{equation}
	where
	\begin{equation}
		\label{eq:HJB-H}
		\begin{aligned}
			H(t,x,p,M,u) = \sup_{a\in A}\Big\{
			&-\dfrac{1}{2}\Tr\big(\sigma(t,x,a)\sigma(t,x,a)^{\top}M\big) \\
			&-\langle\mu(t,x,a),p\rangle
			+ r(t,x,a)\,u - f(t,x,a) \Big\},
		\end{aligned}
	\end{equation}
	\begin{equation}
		\label{eq:HJB-L}
		L(t,x,p,u) = \sup_{b\in B}\Big\{
		\langle\gamma(x,b),p\rangle + k(t,x,b)\,u - g(t,x,b) \Big\}.
 s	\end{equation}
	The data of the problem are as follows:
	\begin{itemize}
		\item $A\subset\R^{N_A}$, $B\subset\R^{N_B}$ are nonempty compact sets;
		\item $\sigma:[0,T]\times\overline{\cO}\times A\to\R^{N\times N_p}$ with $1\le N_p\le N$;
		\item $\mu:[0,T]\times\overline{\cO}\times A\to\R^N$;
		\item $f,r:[0,T]\times\overline{\cO}\times A\to\R$;
		\item $\gamma:\partial\cO\times\mathcal{V}\to\R^N$, where $\mathcal{V}\subseteq\R^{N_B}$ is an open set containing $B$;
		\item $g,k:[0,T]\times\partial\cO\times B\to\R$;
		\item $\Psi:\overline{\cO}\to\R$.
	\end{itemize}
	We denote $\gamma_b(x)=\gamma(x,b)$ for all $(x,b)\in\partial\cO\times B$.
	
	Equations~\eqref{eq:HJB-main}--\eqref{eq:HJB-L} extend the HJB equation with oblique derivative boundary conditions studied by Calzola--Carlini--Dupuis--Silva~\cite{CalzolaEtAl2023}.  The term $r(t,x,a)u$ is the ordinary-time discount, whereas $k(t,x,b)u$ is the oblique Robin term coupled to boundary local time and requires the separate boundary treatment developed below.  When $r\equiv0$ and $k\equiv0$, \eqref{eq:HJB-main}--\eqref{eq:HJB-L} reduce, up to the time reversal $t\mapsto T-t$, to the equation in~\cite{CalzolaEtAl2023}.
	
	
	\subsection{Basic assumptions}
	
	We make the following assumptions throughout the paper.
	
	
	\begin{enumerate}[
		label=\arabic*.,
		ref=\arabic*,
		labelsep=1em,
		leftmargin=*,
		itemindent=0pt
		]
		\item \label{ass:H1}
		$\mathcal{O}\subset\mathbb{R}^N$ is a nonempty bounded domain, and $\partial\mathcal{O}$ is of class $C^3$.
		
		\item \label{ass:H2}
		The functions $\sigma$, $\mu$, $f$, $g$, $r$, $k$ and $\Psi$ are continuous. Moreover, for each $a\in A$, the maps $\sigma(\cdot,\cdot,a)$ and $\mu(\cdot,\cdot,a)$ are Lipschitz continuous, with Lipschitz constants independent of $a\in A$.
		
		\item \label{ass:H3}
		The function $\gamma$ is of class $C^1$, and for every $(x,b)\in\partial\mathcal{O}\times B$,
		\begin{equation}
			|\gamma_b(x)|=1,\qquad \langle n(x),\gamma_b(x)\rangle > 0,
			\label{eq:gamma-cond}
		\end{equation}
		where $n(x)$ denotes the unit outward normal vector of $\overline{\mathcal{O}}$ at $x\in\partial\mathcal{O}$.
		
		\item \label{ass:H4}
		For all $(t,x,a)\in[0,T]\times\overline{\mathcal{O}}\times A$, $r(t,x,a)\ge 0$; and for all $(t,x,b)\in[0,T]\times\partial\mathcal{O}\times B$, $k(t,x,b)\ge 0$.
	\end{enumerate}
	The conditions in Assumptions~\ref{ass:H1}--\ref{ass:H3} on $\cO$, $\sigma$, $\mu$, $f$, $g$, $\Psi$, and $\gamma$ are those used in~\cite{CalzolaEtAl2023}; continuity of the added coefficients $r$ and $k$ is included in Assumption~\ref{ass:H2}.  Assumption~\ref{ass:H4} ensures that the attenuation factors used below do not exceed one; this property is used in the monotonicity and barrier arguments.
	
	
	\subsection{Viscosity solutions and comparison}
	
	The definition of viscosity solutions parallels that of~\cite{CalzolaEtAl2023}, with only the addition of the $ru$ term in the Hamiltonian and the $ku$ term in the boundary operator.
	
	For any bounded function $z:\overline{\cO}_T\to\R$, its upper (resp.\ lower) semicontinuous envelope $z^*$ (resp.\ $z_*$) is defined exactly as in~\cite{CalzolaEtAl2023} and is not repeated here.
	
	\begin{definition}[Viscosity subsolution / supersolution / solution] \upshape
		\label{def:viscosity}
		\hfill\strut
		\begin{enumerate}
			\item[\upshape(i)] An upper semicontinuous function $u_1:\overline{\cO}_T\to\R$ is called a \textbf{viscosity subsolution} of \eqref{eq:HJB-main} if, for every $(t,x)\in\overline{\cO}_T$ and every $\phi\in C^2(\overline{\cO}_T)$ such that $u_1-\phi$ attains a local maximum at $(t,x)$, the following holds:
			\begin{itemize}
				\item If $(t,x)\in\cO_T$,
				\begin{equation}
					-\partial_t\phi(t,x) + H\big(t,x,D\phi(t,x),D^2\phi(t,x),u_1(t,x)\big) \le 0;
					\label{eq:sub-interior}
				\end{equation}
				\item If $(t,x)\in[0,T)\times\partial\cO$,
				\begin{equation}
					\min\Big\{
					-\partial_t\phi(t,x) + H\big(t,x,D\phi(t,x),D^2\phi(t,x),u_1(t,x)\big),\;
					L\big(t,x,D\phi(t,x),u_1(t,x)\big)
					\Big\} \le 0;
					\label{eq:sub-boundary}
				\end{equation}
				\item If $(t,x)\in\{T\}\times\cO$,
				\begin{equation}
					\min\Big\{
					-\partial_t\phi(t,x) + H\big(t,x,D\phi(t,x),D^2\phi(t,x),u_1(t,x)\big),\;
					u_1(t,x)-\Psi(x)
					\Big\} \le 0;
					\label{eq:sub-terminal-interior}
				\end{equation}
				\item If $(t,x)\in\{T\}\times\partial\cO$,
				\begin{equation}
					\begin{split}
						\min\Big\{
						& -\partial_t\phi(t,x) + H\big(t,x,D\phi(t,x),D^2\phi(t,x),u_1(t,x)\big), \\
						& L\big(t,x,D\phi(t,x),u_1(t,x)\big),  u_1(t,x)-\Psi(x)
						\Big\} \le 0.
					\end{split}
					\label{eq:sub-terminal-boundary}
				\end{equation}
			\end{itemize}
			\item[\upshape(ii)] A lower semicontinuous function $u_2$ is called a \textbf{viscosity supersolution} if the inequalities in (i) are reversed: $\le$ is replaced by $\ge$, and $\min$ is replaced by $\max$, with $u_1$ replaced by $u_2$ everywhere.
			
			
			\item[\upshape(iii)] A bounded function $u$ is called a \textbf{viscosity solution} if its upper semicontinuous envelope $u^*$ is a viscosity subsolution and its lower semicontinuous envelope $u_*$ is a viscosity supersolution.
		\end{enumerate}
	\end{definition}
	\begin{remark} \upshape The $\min$ (resp.\ $\max$) form on the terminal layer $\{T\}\times\overline{\cO}$ can be replaced by the pointwise conditions $u_1(T,x)\le\Psi(x)$ (resp.\ $u_2(T,x)\ge\Psi(x)$); see~\cite[Proposition~6]{Bourgoing2008}.
	\end{remark}
	\begin{proposition}[Comparison principle] \upshape\label{thm:wellposed}
		Under Assumptions~\ref{ass:H1}--\ref{ass:H4}, let $v$ be a bounded
		upper semicontinuous viscosity subsolution and let $w$ be a bounded lower
		semicontinuous viscosity supersolution of~\eqref{eq:HJB-main}.  Then
		\[
		v\le w\qquad\text{on }\overline{\cO}_T.
		\]
	\end{proposition}
	\begin{proof}
		Compactness and Assumption~\ref{ass:H3} give
		\[
		\nu:=\min_{(x,b)\in\partial\cO\times B}
		n(x)\cdot\gamma(x,b)>0.
		\]
		Consequently, for every $\lambda\ge0$,
		\[
		L(t,x,p+\lambda n(x),z)
		\ge L(t,x,p,z)+\nu\lambda .
		\]
		Moreover, if $z_1\ge z_2$, then
		\[
		0\le H(t,x,p,M,z_1)-H(t,x,p,M,z_2)
		\le \|r\|_\infty(z_1-z_2),
		\]
		and
		\[
		0\le L(t,x,p,z_1)-L(t,x,p,z_2)
		\le \|k\|_\infty(z_1-z_2).
		\]
		Thus $H$ is continuous, proper, and degenerate elliptic, whereas $L$ is
		continuous, proper, uniformly oblique, and Lipschitz continuous in $p$.
		For example, if $|z|,|w|\le R$ and $d=|t-s|+|x-y|$, uniform moduli
		$\omega_k$ and $\omega_g$ give
		\begin{align*}
		&|L(t,x,p,z)-L(s,y,q,w)|\\
		&\qquad\le |p-q|+C_\gamma|q|\,|x-y|
		+\|k\|_\infty|z-w|+R\omega_k(d)+\omega_g(d).
		\end{align*}
		The uniform Lipschitz continuity of $\sigma$ and $\mu$, together with the
		uniform continuity of the remaining data on their compact domains, gives
		the analogous standard controlled-diffusion modulus for $H$.  Therefore,
		after reversing the
		time variable, the comparison theorem in
		\cite[Theorem~II.1]{Barles1993} applies.  The relaxed terminal formulation
		is equivalent to the pointwise terminal inequalities by
		\cite[Proposition~6]{Bourgoing2008}.
	\end{proof}
	\begin{remark} \upshape
		Existence is not assumed at this stage.  The half-relaxed-limit argument in
		Theorem~\ref{thm:convergence} constructs a viscosity subsolution and a
		viscosity supersolution.  The preceding comparison principle makes them
		coincide and thereby produces the unique continuous viscosity solution.
	\end{remark}
	
	\subsection{Oblique projection near the boundary}
	
	The results in this subsection are taken from~\cite{CalzolaEtAl2023} and are restated here for completeness; they are essential for the scheme construction in Section~3.
	
	A key geometric tool is the existence of a unique oblique projection onto $\partial\cO$ along the direction $\gamma_b$ for points sufficiently close to the boundary.  The following proposition is Proposition~1 of~\cite{CalzolaEtAl2023}, which extends the construction in~\cite[Section~1.2]{Gobet2001} to a control-dependent direction.
	
	\begin{proposition}[Oblique projection] \upshape\label{prop:oblique-proj}
		Under Assumptions~\ref{ass:H1} and~\ref{ass:H3}, there exists $R>0$ such that, for every $x\in\R^N$ satisfying $d(x,\partial\cO)<R$ and every $b\in B$, there exist unique $p^{\gamma_b}(x)\in\partial\cO$ and $d^{\gamma_b}(x)\in\R$ satisfying
		\begin{equation}
			x = p^{\gamma_b}(x) + d^{\gamma_b}(x)\,\gamma_b\!\big(p^{\gamma_b}(x)\big).
			\label{eq:oblique-decomp}
		\end{equation}
		The mappings
		\[
		(x,b)\longmapsto p^{\gamma_b}(x),\qquad (x,b)\longmapsto d^{\gamma_b}(x),
		\]
		called respectively the \textbf{oblique projection} onto $\partial\cO$ and the \textbf{signed algebraic distance} along the direction $\gamma_b$, are both of class $C^1$.
	\end{proposition}
	\begin{proof} \upshape See~\cite{CalzolaEtAl2023}, Proposition~1.\end{proof}
	
	\begin{lemma}[Uniform oblique-distance comparison]
		\upshape\label{lem:oblique-distance}
		After possibly decreasing the radius in
		Proposition~\ref{prop:oblique-proj}, there is a constant $C_\gamma>0$,
		independent of $b\in B$, such that
		\[
		|d^{\gamma_b}(x)|\le C_\gamma\dist(x,\partial\cO)
		\]
		whenever $x$ is in that tubular neighbourhood.  If in addition
		$x\notin\overline\cO$, then $d^{\gamma_b}(x)>0$ and
		\[
		p^{\gamma_b}(x)-d^{\gamma_b}(x)
		\gamma_b\bigl(p^{\gamma_b}(x)\bigr)\in\cO.
		\]
	\end{lemma}
	\begin{proof}
		Let $\rho$ be the signed Euclidean distance to $\partial\cO$, positive
		outside $\cO$ and negative inside.  Since $\partial\cO$ is of class $C^3$,
		$\rho$ is of class $C^2$ in a fixed two-sided tubular neighbourhood and
		$D\rho(p)=n(p)$ on $\partial\cO$.  Compactness and uniform obliqueness give
		\[
		\nu_0:=\min_{p\in\partial\cO,\,b\in B}
		n(p)\cdot\gamma_b(p)>0.
		\]
		On a smaller closed tubular neighbourhood, the joint $C^1$ regularity in
		Proposition~\ref{prop:oblique-proj} gives
		\(
		M:=\sup_{x,b}|D_xd^{\gamma_b}(x)|<\infty.
		\)
		If $\pi(x)$ is the normal projection of $x$ onto $\partial\cO$, then
		$d^{\gamma_b}(\pi(x))=0$.  The segment from $\pi(x)$ to $x$ remains in the
		tubular neighbourhood, and the mean-value theorem therefore yields
		\[
		|d^{\gamma_b}(x)|
		\le M|x-\pi(x)|=M\dist(x,\partial\cO),
		\]
		uniformly in $b$.
		
		It remains to verify the sign and the reflected location.  Taylor's formula,
		uniformly in $(p,b)\in\partial\cO\times B$, gives
		\[
		\rho\bigl(p+s\gamma_b(p)\bigr)
		=s\,n(p)\cdot\gamma_b(p)+O(s^2).
		\]
		Shrink the neighbourhood once more so that the quadratic remainder is at
		most $\nu_0|s|/2$.  For an exterior point write
		$x=p+d\gamma_b(p)$ using Proposition~\ref{prop:oblique-proj}.  Since
		$\rho(x)>0$, the preceding expansion implies $d>0$.  Applying it with
		$s=-d$ gives
		\[
		\rho\bigl(p-d\gamma_b(p)\bigr)
		\le-\nu_0d/2<0,
		\]
		so the reflected point belongs to $\cO$.
	\end{proof}
	
	For any $\varepsilon\ge0$, we introduce the following notation :
	\begin{equation}
		\begin{aligned}
			D_\varepsilon    &:= \{x\in\overline{\cO}\mid d(x,\partial\cO)>\varepsilon\},\\[2pt]
			\partial D_\varepsilon &:= \{x\in\overline{\cO}\mid d(x,\partial\cO)=\varepsilon\},\\[2pt]
			L_\varepsilon    &:= \{x\in\overline{\cO}\mid d(x,\partial\cO)\le\varepsilon\}.
		\end{aligned}
		\label{eq:domain-decomp}
	\end{equation}
	
	\begin{lemma} [Regularity of the distance function]. \upshape\label{lem:dist-reg}
		Under Assumption~\ref{ass:H1}, the following hold:
		\begin{enumerate}
			\item[\upshape(i)] There exists $\eta>0$ such that on $L_\eta$, the normal projection $p_{\partial\cO}$ is well defined and of class $C^1$.
			\item[\upshape(ii)] The distance function $L_\eta\ni x\mapsto d(x,\partial\cO)$ is of class $C^3$, and
			\[
			D d(\cdot,\partial\cO)(x) = -\,n\!\big(p_{\partial\cO}(x)\big).
			\]
			\item[\upshape(iii)] For every $\delta\in[0,\eta]$, $\partial D_\delta$ is of class $C^3$, and its unit outward normal is $n_\delta(x)=n(p_{\partial\cO}(x))$.
			\item[\upshape(iv)] For any $x\in L_\delta$ ($\delta\in[0,\eta]$), the point $p = p_{\partial\cO}(x) - \delta\,n(p_{\partial\cO}(x))$ is the projection of $x$ onto $\partial D_\delta$.
			\item[\upshape(v)] For every $\delta\in[0,\eta]$, the function $x\mapsto d(x,\partial D_\delta)$ is of class $C^3$ on $L_\delta$, and
			\begin{equation}
				d(x,\partial\cO) + d(x,\partial D_\delta) = \delta,\qquad \forall\,x\in L_\delta.
				\label{eq:dist-add}
			\end{equation}
		\end{enumerate}
	\end{lemma}
	\begin{proof}
		\upshape See~\cite[Lemma~1]{CalzolaEtAl2023}; parts (i) and (ii) are taken from~\cite[Lemma~14.16]{GilbargTrudinger2001}.
	\end{proof}
	
	
	% ============================================================
	% ============================================================
	\section{Probabilistic construction of the fully discrete scheme}
	\label{sec:scheme}
	
	\subsection{Spatial grid and \texorpdfstring{$\mathbb P_1$}{P1} interpolation}
	
	Let $\cO_h$ be the polyhedral approximation of $\cO$ and let $\mathbb T_h$
	be a fitted, shape-regular simplicial triangulation constructed as
	in~\cite[Section~3]{CalzolaEtAl2023}.  Its spatial mesh size is
	\[
	h:=\max_{\mathsf T\in\mathbb T_h}\operatorname{diam}(\mathsf T)
	\]
	and tends to zero along the mesh family.  Write
	\[
	\cG_h=\{x_i:i\in\mathcal I_h\}\subset\overline\cO,\qquad
	\mathcal I_h=\mathcal I_h^\circ\mathbin{\dot\cup}\mathcal I_h^\partial,
	\]
	where $x_i\in\cO$ for $i\in\mathcal I_h^\circ$ and
	$x_i\in\partial\cO$ for $i\in\mathcal I_h^\partial$.  As in the cited
	construction, $\widehat{\mathbb T}_h$ denotes the corresponding exact
	triangulation of $\overline\cO$ with the same vertices and with curved
	boundary elements.
	
	Let $\{\psi_j\}_{j\in\mathcal I_h}$ be the continuous piecewise-linear
	($\mathbb P_1$) nodal basis on $\mathbb T_h$, and let $p_h$ be the
	elementwise map from $\widehat{\mathbb T}_h$ to $\mathbb T_h$ used
	in~\cite[Section~3]{CalzolaEtAl2023}.  For a grid vector
	$V=(V_j)_{j\in\mathcal I_h}$ define
	\begin{equation}
		I_h[V](x):=\sum_{j\in\mathcal I_h}\psi_j(p_h(x))V_j,
		\qquad x\in\overline\cO.
		\label{eq:interpolation}
	\end{equation}
	This $\mathbb P_1$ interpolant is the only interpolation operator used below.
	Its barycentric weights are nonnegative and sum to one; hence $I_h$ is
	positivity preserving.  Moreover, the standard estimate proved for this
	mesh construction in~\cite[Lemma~2]{CalzolaEtAl2023} gives, for every
	$\phi\in C^2(\overline\cO)$,
	\begin{equation}
		\bigl\|I_h[\phi|_{\cG_h}]-\phi\bigr\|_{L^\infty(\overline\cO)}
		\le C_\phi h^2 .
		\label{eq:interp-error}
	\end{equation}
	
	Choose $N_T\ge1$ and set
	\[
	\dt=\frac{T}{N_T},\qquad t_n=n\dt,\qquad n=0,\ldots,N_T.
	\]
	\subsection{Reflected Feynman--Kac representation}
	
	We first suppress the controls in order to show where every term of the update comes from.  Consider
	\begin{equation}
		\begin{cases}
			-u_t-\dfrac12\Tr(\bar\sigma\bar\sigma^\top D^2u)-\bar\mu\cdot Du+\bar r u-\bar f=0,&(t,x)\in \cO_T,\\
			\bar\gamma\cdot Du+\bar k u=\bar g,&(t,x)\in[0,T)\times\partial\cO,\\
			u(T,x)=\Psi(x),&x\in\overline\cO.
		\end{cases}
		\label{eq:linear-problem}
	\end{equation}
	For this derivation, assume temporarily that a reflected semimartingale pair $(X,L)$ solving the equation below exists, that $u\in C^{1,2}([0,T]\times\overline\cO)$, and that the coefficients and $u$ have the boundedness and integrability needed for It\^{o}'s formula and conditional expectation.  These are auxiliary classical hypotheses only; the convergence proof below neither invokes the representation nor assumes this regularity.  Let
	\[
	dX_s=\bar\mu(s,X_s)\,ds+\bar\sigma(s,X_s)\,dW_s-\bar\gamma(X_s)\,dL_s,
	\]
	where $L$ is nondecreasing and increases only on $\partial\cO$.  The minus sign is consistent with $\bar\gamma$ pointing outward.  For $s\ge t_n$, define
	\[
	A_s=\exp\!\left(-\int_{t_n}^s\bar r(q,X_q)\,dq-\int_{t_n}^s\bar k(q,X_q)\,dL_q\right).
	\]
	It\^{o}'s formula gives the one-step conditional representation
	\begin{equation}
		\begin{aligned}
			u(t_n,x_i)=\E\Bigg[&A_{t_{n+1}}u(t_{n+1},X_{t_{n+1}})
			+\int_{t_n}^{t_{n+1}}A_s\bar f(s,X_s)\,ds\\
			&+\int_{t_n}^{t_{n+1}}A_s\bar g(s,X_s)\,dL_s
			\,\Big|\,X_{t_n}=x_i\Bigg].
		\end{aligned}
		\label{eq:FK}
	\end{equation}
	This identity, rather than a characteristic interpolation formula, is the starting point of the discrete method.
	
	\subsection{Equal-probability discretization of the linear representation}
	
	Set $m:=N_p$, the Brownian dimension.  Let $\xi\in\R^m$ have independent Rademacher components:
	\[
	\mathbb P(\xi_\ell=1)=\mathbb P(\xi_\ell=-1)=\frac12,
	\qquad \ell=1,\ldots,m.
	\]
	Its $P=2^m$ possible values are denoted by $\xi^1,\ldots,\xi^P$.  Every value occurs with probability $1/P$, and
	\begin{equation}
		\frac1P\sum_{p=1}^P\xi^p=0,
		\qquad
		\frac1P\sum_{p=1}^P\xi^p(\xi^p)^\top=I_m.
		\label{eq:cubature}
	\end{equation}
	These equal probabilities are part of the definition of the scheme and are not free numerical parameters.
	
	We first discretize the linear problem~\eqref{eq:linear-problem}.  Write
	$\overline\gamma=\gamma(\cdot,\overline b)$ for its fixed boundary
	parameter and abbreviate
	$p^{\overline\gamma}=p^{\gamma_{\overline b}}$ and
	$d^{\overline\gamma}=d^{\gamma_{\overline b}}$.  We retain the trial-point notation used later for the HJB
	scheme.  Its $p$th weak-Euler trial point is
	\begin{equation}
		X_{n,i}^{*,p}=x_i+\dt\,\overline\mu(t_n,x_i)
		+\sqrt{\dt}\,\overline\sigma(t_n,x_i)\xi^p.
		\label{eq:linear-trial}
	\end{equation}
	If $X_{n,i}^{*,p}\in\overline\cO$, put
	\[
	\widetilde X_{n,i}^p=X_{n,i}^{*,p},
	\qquad
	\widetilde d_{n,i}^p=\widetilde k_{n,i}^p=\widetilde g_{n,i}^p=0.
	\]
	If $X_{n,i}^{*,p}\notin\overline\cO$, let
	\begin{align}
		x_{n,i}^{\overline b,p}&=p^{\overline\gamma}(X_{n,i}^{*,p}),
		&\widetilde d_{n,i}^p&=2d^{\overline\gamma}(X_{n,i}^{*,p}),
		\label{eq:linear-projection}\\
		\widetilde X_{n,i}^p
		&=X_{n,i}^{*,p}-\widetilde d_{n,i}^p
		\overline\gamma(x_{n,i}^{\overline b,p}),
		\label{eq:linear-mirror}\\
		\widetilde k_{n,i}^p&=\overline k(t_n,x_{n,i}^{\overline b,p}),
		&\widetilde g_{n,i}^p&=\overline g(t_n,x_{n,i}^{\overline b,p}).
		\label{eq:linear-boundary-data}
	\end{align}
	The returned point $\widetilde X_{n,i}^p$ is the mirror image of the trial
	point about its oblique projection, and the round-trip distance
	$\widetilde d_{n,i}^p$ is used as the discrete increment of the regulator
	$L$ in~\eqref{eq:FK}.
	
	We now approximate the three terms in~\eqref{eq:FK} separately.  First freeze the ordinary-time coefficient at $(t_n,x_i)$.  The scalar discount is approximated positively by
	\begin{equation}
		e^{-\overline r(t_n,x_i)\dt}
		=\frac{1}{1+\overline r(t_n,x_i)\dt}+O(\dt^2).
		\label{eq:ordinary-discount}
	\end{equation}
	Freezing the boundary coefficient on branch $p$ then gives the continuation approximation
	\begin{equation}
		\frac{1}{1+\overline r(t_n,x_i)\dt}\,
		\frac1P\sum_{p=1}^P
		e^{-\widetilde k_{n,i}^p\widetilde d_{n,i}^p}
		u(t_{n+1},\widetilde X_{n,i}^p).
		\label{eq:linear-continuation}
	\end{equation}
	Second, the ordinary-time source is approximated by the left-endpoint rule,
	\begin{equation}
		\E\!\left[\int_{t_n}^{t_{n+1}}A_s\overline f(s,X_s)\,ds\right]
		\approx\dt\,\overline f(t_n,x_i).
		\label{eq:linear-volume-source}
	\end{equation}
	Third, on a branch with local-time increment $D$, freezing $K$ and $G$ gives
	\[
	\int_0^D e^{-K\ell}G\,d\ell.
	\]
	The midpoint rule yields
	\begin{equation}
		\int_0^D e^{-K\ell}G\,d\ell
		=D e^{-KD/2}G+O(D^3).
		\label{eq:midpoint-local-time}
	\end{equation}
	Since $\widetilde d_{n,i}^p=O(\sqrt\dt)$, the boundary-source contribution becomes
	\begin{equation}
		\frac{1}{1+\overline r(t_n,x_i)\dt}\,
		\frac1P\sum_{p=1}^P
		\widetilde d_{n,i}^p e^{-\widetilde k_{n,i}^p\widetilde d_{n,i}^p/2}
		\widetilde g_{n,i}^p,
		\label{eq:linear-boundary-source}
	\end{equation}
	up to the local $O(\dt^{3/2})$ quadrature error used later in the consistency proof.
	
	Replacing the next-time solution by a grid vector $V$ and using the $\mathbb P_1$ interpolant gives the linear interior operator
	\begin{equation}
		\begin{aligned}
			\overline{\cS}_{n,i}[V]
			=\frac{1}{1+\overline r(t_n,x_i)\dt}\frac1P\sum_{p=1}^P
			\Big[&e^{-\widetilde k_{n,i}^p\widetilde d_{n,i}^p}
			I_h[V](\widetilde X_{n,i}^p)\\
			&+\widetilde d_{n,i}^p
			e^{-\widetilde k_{n,i}^p\widetilde d_{n,i}^p/2}
			\widetilde g_{n,i}^p\Big]
			+\dt\,\overline f(t_n,x_i).
		\end{aligned}
		\label{eq:linear-scheme}
	\end{equation}
	For the boundary equation choose a purely spatial offset $\ell_h$, independent
	of $\dt$, such that $c_0h\le\ell_h\le c_1h$ for two fixed constants
	$0<c_0\le c_1$.  The complete linear method is
	\begin{equation}
		\begin{cases}
			U_i^n=\overline{\cS}_{n,i}[U^{n+1}],
			&n=N_T-1,\ldots,0,\quad i\in\mathcal I_h^\circ,\\[6pt]
			\displaystyle
			\frac{U_i^n-I_h[U^n](x_i-\ell_h\overline\gamma(x_i))}{\ell_h}
			+\overline k(t_n,x_i)U_i^n-\overline g(t_n,x_i)=0,
			&n=N_T-1,\ldots,0,\quad i\in\mathcal I_h^\partial,\\[9pt]
			U_i^{N_T}=\Psi(x_i),&i\in\mathcal I_h .
		\end{cases}
		\label{eq:complete-linear-scheme}
	\end{equation}
	Thus one first computes the interior values from $U^{n+1}$ and then solves
	the same-level boundary equations.  These equations can be implicit because
	the simplex containing $x_i-\ell_h\overline\gamma(x_i)$ may have boundary
	vertices.  Their coefficient matrix is a sparse nonsingular $M$-matrix by
	the $\mathbb P_1$ weight argument stated after the general scheme below.
	The coefficient $\overline k$ also acts on every exiting interior branch
	through local-time attenuation, while every branch probability remains
	exactly $1/P$.
	\subsection{Extension to the controlled HJB problem}
	
	We now return to~\eqref{eq:HJB-main}.  For one fixed evaluation of the interior
	update at $(t_n,x_i)$, choose $a\in A$ and $b\in B$ and freeze this pair over
	the corresponding one-step transition.  The same $b$ is used for all $P$
	equally probable branches emanating from $x_i$; a different boundary control
	is not selected on each branch.  After the $1/P$-weighted branch average has
	been formed, the infimum over $(a,b)$ is taken at that node.  Controls may
	therefore be selected anew at different nodes and time levels.  Using the
	notation of the original manuscript, the $p$th trial point is
	\begin{equation}
		X_{n,i}^{*,p}(a)=x_i+\dt\,\mu(t_n,x_i,a)
		+\sqrt{\dt}\,\sigma(t_n,x_i,a)\xi^p.
		\label{eq:trial}
	\end{equation}
	If $X_{n,i}^{*,p}(a)\in\overline\cO$, set
	\[
	\widetilde X_{n,i}^p(a,b)=X_{n,i}^{*,p}(a),\qquad
	\widetilde d_{n,i}^p(a,b)=\widetilde k_{n,i}^p(a,b)=
	\widetilde g_{n,i}^p(a,b)=0.
	\]
	If $X_{n,i}^{*,p}(a)\notin\overline\cO$, set
	\begin{align}
		x_{n,i}^{b,p}(a)&=p^{\gamma_b}(X_{n,i}^{*,p}(a)),\\
		\widetilde d_{n,i}^p(a,b)&=2d^{\gamma_b}(X_{n,i}^{*,p}(a)),\\
		\widetilde X_{n,i}^p(a,b)&=X_{n,i}^{*,p}(a)
		-\widetilde d_{n,i}^p(a,b)\gamma_b(x_{n,i}^{b,p}(a)),\\
		\widetilde k_{n,i}^p(a,b)&=k(t_n,x_{n,i}^{b,p}(a),b),\\
		\widetilde g_{n,i}^p(a,b)&=g(t_n,x_{n,i}^{b,p}(a),b).
		\label{eq:branch-data}
	\end{align}
	Proposition~\ref{prop:oblique-proj} makes these quantities well defined for sufficiently small $\dt$.  Uniformly in all nodes and controls,
	\begin{equation}
		|X_{n,i}^{*,p}(a)-x_i|+\widetilde d_{n,i}^p(a,b)\le C\sqrt{\dt}.
		\label{eq:step-bound}
	\end{equation}
	Indeed, continuity on the compact coefficient set and finiteness of the
	Rademacher family give, for $\dt\le1$,
	\[
	|X_{n,i}^{*,p}(a)-x_i|
	\le\dt\|\mu\|_\infty
	+\sqrt\dt\|\sigma\|_\infty\max_{1\le q\le P}|\xi^q|
	\le C_0\sqrt\dt .
	\]
	If the trial point is exterior, the segment joining $x_i\in\cO$ to
	$X_{n,i}^{*,p}(a)$ meets $\partial\cO$, and hence
	\[
	\dist(X_{n,i}^{*,p}(a),\partial\cO)
	\le|X_{n,i}^{*,p}(a)-x_i|.
	\]
	For sufficiently small $\dt$, Lemma~\ref{lem:oblique-distance} applies and
	gives
	\[
	0<\widetilde d_{n,i}^p(a,b)
	=2d^{\gamma_b}(X_{n,i}^{*,p}(a))
	\le2C_\gamma|X_{n,i}^{*,p}(a)-x_i|
	\le C\sqrt\dt.
	\]
	The same lemma shows that $\widetilde X_{n,i}^p(a,b)\in\cO$.  For a
	non-exiting branch $\widetilde d_{n,i}^p=0$ and
	$\widetilde X_{n,i}^p=X_{n,i}^{*,p}\in\overline\cO$.  Thus
	\eqref{eq:step-bound} holds and every interpolation value below is well
	defined.
	
	Applying the already derived linear recurrence with the frozen coefficients gives, for a grid vector $V$, the following operator.  In its defining sum only, we suppress the common arguments $(n,i,a,b)$:
	\begin{equation}
		\begin{aligned}
			\cS_{n,i}^{a,b}[V]
			=\frac{1}{1+r(t_n,x_i,a)\dt}\frac1P\sum_{p=1}^P
			\Big[e^{-\widetilde k_p\widetilde d_p}I_h[V](\widetilde X_p)
			+\widetilde d_p e^{-\widetilde k_p\widetilde d_p/2}\widetilde g_p\Big]
			+\dt f(t_n,x_i,a).
		\end{aligned}
		\label{eq:fixed-operator}
	\end{equation}
	For $i\in\mathcal I_h^\circ$, the HJB operator minimizes the fixed-control conditional expectation:
	\begin{equation}
		S_{n,i}[V]=\inf_{(a,b)\in A\times B}\cS_{n,i}^{a,b}[V].
		\label{eq:HJB-operator}
	\end{equation}
	
	For $i\in\mathcal I_h^\partial$ and $b\in B$, set
	\[
	y_{i,b}:=x_i-\ell_h\gamma(x_i,b).
	\]
	Uniform obliqueness implies $y_{i,b}\in\cO$ for all sufficiently small $h$,
	uniformly in $i$ and $b$.  Set
	\[
	\omega_{ij}^b:=\psi_j(p_h(y_{i,b})).
	\]
	The fixed-control and controlled boundary updates are
	\begin{align}
		\cS_{n,i}^{\partial,b}[V]
		&:=\frac{I_h[V](y_{i,b})+\ell_h\,g(t_n,x_i,b)}
		{1+\ell_h\,k(t_n,x_i,b)},
		\label{eq:fixed-boundary-map}\\
		S_{n,i}^{\partial}[V]&:=\inf_{b\in B}\cS_{n,i}^{\partial,b}[V].
		\label{eq:controlled-boundary-map}
	\end{align}
	Since $1+\ell_h\,k(t_n,x_i,b)>0$ and $B$ is compact,
	$V_i=S_{n,i}^{\partial}[V]$ is equivalent to
	\begin{equation}
		\sup_{b\in B}
		\left\{
		\frac{V_i-I_h[V](x_i-\ell_h\gamma(x_i,b))}{\ell_h}
		+k(t_n,x_i,b)V_i-g(t_n,x_i,b)
		\right\}=0.
		\label{eq:discrete-boundary-residual}
	\end{equation}
	
	The complete nonlinear method is therefore
	\begin{equation}
		\begin{cases}
			U_i^n=S_{n,i}[U^{n+1}],
			&n=N_T-1,\ldots,0,\quad i\in\mathcal I_h^\circ,\\[7pt]
			\begin{aligned}
				\sup_{b\in B}\biggl\{&
				\frac{U_i^n-I_h[U^n](x_i-\ell_h\gamma(x_i,b))}{\ell_h}\\[-2pt]
				&+k(t_n,x_i,b)U_i^n-g(t_n,x_i,b)
				\biggr\}=0,
			\end{aligned}
			&n=N_T-1,\ldots,0,\quad i\in\mathcal I_h^\partial,\\[10pt]
			U_i^{N_T}=\Psi(x_i),&i\in\mathcal I_h .
		\end{cases}
		\label{eq:scheme}
	\end{equation}
	
	\begin{lemma}[Uniform interior interpolation weight]
		\label{lem:uniform-interior-weight}
		There exist $h_0>0$ and $\theta\in(0,1)$, independent of $h$, $i$, and $b$,
		such that, for $0<h\le h_0$,
		\begin{equation}
			\sum_{j\in\mathcal I_h^\circ}\omega_{ij}^b\ge\theta .
			\label{eq:boundary-row-margin}
		\end{equation}
	\end{lemma}
	\begin{proof}
		Let
		\[
		\nu:=\min_{(x,b)\in\partial\cO\times B}
		n(x)\cdot\gamma(x,b)>0,
		\]
		and choose $\varrho\in C^2(\overline\cO)$ that agrees with
		$\dist(\cdot,\partial\cO)$ in a fixed inner tubular neighbourhood.  Since
		$D\varrho=-n$ on $\partial\cO$, Taylor's formula and
		$c_0h\le\ell_h\le c_1h$ give, uniformly in $i$ and $b$,
		\[
		\varrho(y_{i,b})
		=\ell_h n(x_i)\cdot\gamma(x_i,b)+O(\ell_h^2)
		\ge c_0\nu h-Ch^2.
		\]
		The interpolation estimate~\eqref{eq:interp-error} therefore implies, for
		all sufficiently small $h$,
		\[
		I_h[\varrho|_{\cG_h}](y_{i,b})\ge\frac{c_0\nu}{2}h.
		\]
		
		Let $\mathsf T\in\mathbb T_h$ be an active simplex for $p_h(y_{i,b})$.
		If $\mathsf T$ has no boundary vertex, every nonzero barycentric weight is
		attached to an interior node and the sum in~\eqref{eq:boundary-row-margin}
		equals one.  Otherwise choose a boundary vertex
		$x_q\in\mathsf T\cap\partial\cO$.  Every vertex $x_j$ of $\mathsf T$
		satisfies
		\[
		\dist(x_j,\partial\cO)\le|x_j-x_q|
		\le\operatorname{diam}(\mathsf T)\le h.
		\]
		For small $h$, $\varrho(x_j)=\dist(x_j,\partial\cO)$ at all these vertices,
		and $\varrho$ vanishes at the boundary vertices.  Locality and nonnegativity
		of the $\mathbb P_1$ weights now give
		\[
		I_h[\varrho|_{\cG_h}](y_{i,b})
		=\sum_j\omega_{ij}^b\varrho(x_j)
		\le h\sum_{j\in\mathcal I_h^\circ}\omega_{ij}^b.
		\]
		The last two estimates prove~\eqref{eq:boundary-row-margin}, for example
		with $\theta=\min\{1/2,c_0\nu/2\}$.  This is precisely where the lower
		offset bound $\ell_h\ge c_0h$ is used.
	\end{proof}
	
	After the interior values at time level $t_n$ have been computed, regard the
	right-hand side of~\eqref{eq:controlled-boundary-map} as a map $T_n$ of the
	boundary vector.  Lemma~\ref{lem:uniform-interior-weight}, $k\ge0$, and the
	partition-of-unity property give directly
	\[
	\|T_nz-T_n\widetilde z\|_\infty
	\le\sup_{i,b}
	\frac{\sum_{j\in\mathcal I_h^\partial}\omega_{ij}^b}
	{1+\ell_hk(t_n,x_i,b)}
	\|z-\widetilde z\|_\infty
	\le(1-\theta)\|z-\widetilde z\|_\infty.
	\]
	Hence the boundary values are obtained by the contraction iteration
	$z^{(q+1)}=T_nz^{(q)}$.  In the linear case the same row-margin calculation
	shows directly that the boundary coefficient matrix is a sparse nonsingular
	$M$-matrix.  Indeed, writing $k_i=k(t_n,x_i)$ for fixed linear data, its
	boundary block has entries
	\[
	M_{ii}=1+\ell_h k_i-\omega_{ii},\qquad
	M_{ij}=-\omega_{ij}\quad(i\ne j),
	\]
	and its strict row margin is
	\[
	M_{ii}-\sum_{j\ne i}|M_{ij}|
	=\ell_h k_i+\sum_{j\in\mathcal I_h^\circ}\omega_{ij}
	\ge\theta.
	\]
	\begin{remark}[Coupling of the two Robin weights]\label{rem:coupled-weights}
		Let $x_b\in\partial\cO$,
		$X^*=x_b+(\widetilde d/2)\gamma$, and
		$\widetilde X=x_b-(\widetilde d/2)\gamma$.  For a smooth test function,
		Taylor expansion gives
		\[
		e^{-k\widetilde d}\phi(\widetilde X)
		+\widetilde d e^{-k\widetilde d/2}g-\phi(X^*)
		=-\widetilde d\bigl(\gamma\cdot D\phi+k\phi-g\bigr)(x_b)
		+\frac{k\widetilde d^2}{2}
		\bigl(\gamma\cdot D\phi+k\phi-g\bigr)(x_b)
		+O(\widetilde d^3).
		\]
		Thus $e^{-k\widetilde d}$ is not a modification of the branch probability $1/P$; it is
		the local-time attenuation of the continuation value.  The midpoint factor in
		the boundary source makes the entire quadratic term a multiple of the same
		Robin residual.  This elementary identity is used explicitly in the
		near-boundary branch expansion below.
	\end{remark}
	
	\begin{remark}[Polygonal implementations]
		The convergence theorem below is proved for a $C^3$ boundary, where the control-uniform oblique projection is available.  In the disk experiments the branch projection and reflection use the exact circle, while interpolation is performed on a fitted regular polygon after an $O(h^2)$ radial cap map.  On a rectangle, a numerical branch can be reflected face by face using the constant outward normal of each face.  Mixed Robin--Dirichlet stopping requires a separate first-hit rule and is reported only as a numerical extension; it is not used as evidence for the smooth-domain convergence theorem.
	\end{remark}
	
	The role of the Robin term is now unambiguous.  When $\widetilde k\widetilde d>0$, the coefficient of the unknown in an exiting branch is $e^{-\widetilde k\widetilde d}<1$.  No change of the prescribed source alone can reproduce this coefficient for every continuation vector.  This observation distinguishes the local-time attenuation from an additive modification of an oblique-derivative boundary update, without changing the equal probability assigned to the branch.
	
	\section{Properties of the fully discrete scheme}
	\label{sec:properties}
	
	For $(t,x,p,M,z,a)\in\overline{\cO}_T\times\R^N\times\R^{N\times N}\times\R\times A$
	and $(t,x,p,z,b)\in[0,T]\times\partial\cO\times\R^N\times\R\times B$, define
	the control-resolved operators
	\begin{equation}
		\begin{aligned}
			\mathcal H(t,x,p,M,z,a)
			&:=-\frac12\Tr\!\left(\sigma\sigma^\top(t,x,a)M\right)
			-\mu(t,x,a)\cdot p+r(t,x,a)z-f(t,x,a),\\
			\mathcal L(t,x,p,z,b)
			&:=\gamma(x,b)\cdot p+k(t,x,b)z-g(t,x,b).
		\end{aligned}
		\label{eq:calH-calL}
	\end{equation}
	Thus $H=\sup_{a\in A}\mathcal H$ and $L=\sup_{b\in B}\mathcal L$.  As in the
	original consistency calculation, set
	\begin{equation}
		\mathcal R(t,x,p,z,b)
		:=\frac12 k(t,x,b)^2z-\frac12 k(t,x,b)g(t,x,b)
		+\frac12 k(t,x,b)\gamma(x,b)\cdot p.
		\label{eq:calR}
	\end{equation}
	In particular,
	\begin{equation}
		\mathcal R(t,x,p,z,b)=\frac12 k(t,x,b)\mathcal L(t,x,p,z,b).
		\label{eq:R-equals-kL}
	\end{equation}
	For an exiting branch let
	$x_b^p:=x_{n,i}^{b,p}(a)=p^{\gamma_b}(X_{n,i}^{*,p}(a))\in\partial\cO$ and define
	\begin{equation}
		\widetilde{\mathcal L}_{n,i}^p(q,v,a,b):=
		\begin{cases}
			0,&X_{n,i}^{*,p}(a)\in\overline\cO,\\[3pt]
			\mathcal L(t_n,x_b^p,q,v,b),&X_{n,i}^{*,p}(a)\notin\overline\cO,
		\end{cases}
		\label{eq:tildeL}
	\end{equation}
	and define $\widetilde{\mathcal R}_{n,i}^p$ in the same way with
	$\mathcal L$ replaced by $\mathcal R$.  On a non-exiting branch the displayed
	arguments involving $x_b^p$ are ignored because both truncated operators are
	defined to be zero.
	
	For a smooth function $\phi$, we use the original shorthand
	\[
	\bigl[-\partial_t\phi+\mathcal H\bigr]
	(t,x,D\phi,D^2\phi,\phi,a)
	:=-\partial_t\phi(t,x)
	+\mathcal H(t,x,D\phi(t,x),D^2\phi(t,x),\phi(t,x),a).
	\]
	
	\subsection{Monotonicity of the parabolic scheme}
	
	\begin{proposition}[Parabolic monotonicity]\label{prop:monotone}
		Under Assumptions~\ref{ass:H1}--\ref{ass:H4}, let $V\le W$ be two grid
		functions.  For every $n=0,\ldots,N_T-1$ the following assertions hold:
		\begin{align}
			\cS_{n,i}^{a,b}[V]&\le\cS_{n,i}^{a,b}[W],
			&S_{n,i}[V]&\le S_{n,i}[W],
			&&i\in\mathcal I_h^\circ, (a,b)\in A\times B,
			\label{eq:para-monotone-interior}\\
			\cS_{n,i}^{\partial,b}[V]&\le\cS_{n,i}^{\partial,b}[W],
			&S_{n,i}^{\partial}[V]&\le S_{n,i}^{\partial}[W],
			&&i\in\mathcal I_h^\partial, b\in B.
			\label{eq:para-monotone-boundary}
		\end{align}
		Consequently, one complete backward time-layer update is order preserving.
		Moreover, for every constant $c\ge0$,
		\begin{equation}
			S_{n,i}[V+c]\le S_{n,i}[V]+c,\qquad
			S_{n,i}^{\partial}[V+c]\le S_{n,i}^{\partial}[V]+c.
			\label{eq:para-shift}
		\end{equation}
		If $f,g,$ and $\Psi$ are nonnegative, then the discrete solution is
		nonnegative.
	\end{proposition}
	
	\begin{proof}
		All assertions follow directly from~\eqref{eq:scheme}: the interpolation and branch coefficients are nonnegative, their continuation coefficients have total mass at most one, infima preserve order, and the same-level boundary equation is solved by an order-preserving contraction.
	\end{proof}
	
	For completeness, the constant in the shift estimate is explicit.  Put
	\[
	C_{rk}:=\max\{\|r\|_\infty,\|k\|_\infty\}.
	\]
	For every
	$i\in\mathcal I_h^\circ$, fixed $(a,b)$, and constant $c$, adding $c$ to
	$V$ multiplies it by
	\[
	\alpha_{n,i}^{a,b}:=
	\frac{P^{-1}\sum_{p=1}^P
		e^{-\widetilde k_{n,i}^p(a,b)\widetilde d_{n,i}^p(a,b)}}
	{1+r(t_n,x_i,a)\dt}\in[0,1].
	\]
	Since $1-e^{-z}\le z$,
	\begin{equation}
		\left|\cS_{n,i}^{a,b}[V+c]-\cS_{n,i}^{a,b}[V]-c\right|
		=|c|(1-\alpha_{n,i}^{a,b})
		\le C_{rk}|c|\left(\dt+\frac1P\sum_{p=1}^P
		\widetilde d_{n,i}^p(a,b)\right).
		\label{eq:shift-defect}
	\end{equation}
	
	\subsection{Consistency of the parabolic scheme}
	
	\begin{proposition}[Parabolic consistency]\label{prop:nodewise-consistency}
		Let $\phi$ be the restriction to $[0,T]\times\overline\cO$ of a $C^3$
		function defined on a fixed neighbourhood, with all derivatives used below
		bounded there, and put $\phi_{n,i}:=\phi(t_n,x_i)$.  For sufficiently small
		$\Delta t$ and $h$, the following statements hold uniformly in the indicated nodes
		and controls.
		
		\smallskip
		\noindent\textup{(I) Interior nodes.}
		For $i\in\mathcal I_h^\circ$, $a\in A$, and $b\in B$,
		\begin{align}
			&\cS_{n,i}^{a,b}\bigl[\phi(t_{n+1},\cdot)|_{\cG_h}\bigr]-\phi_{n,i}
			\notag\\
			&\quad=-\Delta t\,
			\bigl[-\partial_t\phi+\mathcal H\bigr]
			\bigl(t_n,x_i,D\phi_{n,i},D^2\phi_{n,i},\phi_{n,i},a\bigr)
			\notag\\
			&\qquad-\frac1P\sum_{p=1}^P\widetilde d_{n,i}^p(a,b)
			\widetilde{\mathcal L}_{n,i}^p
			\bigl(D\phi(t_n,x_b^p),\phi(t_n,x_b^p),a,b\bigr)
			\notag\\
			&\qquad+\frac1P\sum_{p=1}^P\bigl(\widetilde d_{n,i}^p(a,b)\bigr)^2
			\widetilde{\mathcal R}_{n,i}^p
			\bigl(D\phi(t_n,x_b^p),\phi(t_n,x_b^p),a,b\bigr)
			+O(\Delta t^{3/2}+h^2).
			\label{eq:reflected-expansion-fixed}
		\end{align}
		Consequently,
		\begin{align}
			&S_{n,i}\bigl[\phi(t_{n+1},\cdot)|_{\cG_h}\bigr]-\phi_{n,i}
			\notag\\
			&\quad=-\sup_{a\in A,\,b\in B}\Biggl\{
			\Delta t\,
			\bigl[-\partial_t\phi+\mathcal H\bigr]
			\bigl(t_n,x_i,D\phi_{n,i},D^2\phi_{n,i},\phi_{n,i},a\bigr)
			\notag\\
			&\hspace{25mm}+\frac1P\sum_{p=1}^P\widetilde d_{n,i}^p(a,b)
			\widetilde{\mathcal L}_{n,i}^p
			\bigl(D\phi(t_n,x_b^p),\phi(t_n,x_b^p),a,b\bigr)
			\notag\\
			&\hspace{25mm}-\frac1P\sum_{p=1}^P\bigl(\widetilde d_{n,i}^p(a,b)\bigr)^2
			\widetilde{\mathcal R}_{n,i}^p
			\bigl(D\phi(t_n,x_b^p),\phi(t_n,x_b^p),a,b\bigr)
			\Biggr\}
			+O(\Delta t^{3/2}+h^2).
			\label{eq:reflected-expansion-controlled}
		\end{align}
		
		\smallskip
		\noindent\textup{(II) Boundary nodes.}
		For $i\in\mathcal I_h^\partial$,
		\begin{equation}
			S_{n,i}^{\partial}\bigl[\phi(t_n,\cdot)|_{\cG_h}\bigr]-\phi_{n,i}
			=-\ell_h\sup_{b\in B}
			\mathcal L(t_n,x_i,D\phi_{n,i},\phi_{n,i},b)+O(h^2).
			\label{eq:boundary-controlled-consistency}
		\end{equation}
		The boundary formula is entirely at time $t_n$; it contains neither a time
		derivative nor the volume Hamiltonian.
	\end{proposition}
	
	\begin{proof}
		Throughout the proof, $C>0$ denotes a generic constant independent of the
		discretization parameters.  Fix $n$, $i\in\mathcal I_h^\circ$, $a\in A$,
		$b\in B$, and put
		\[
		\mu=\mu(t_n,x_i,a),\quad \sigma=\sigma(t_n,x_i,a),\quad
		r=r(t_n,x_i,a),\quad f=f(t_n,x_i,a).
		\]
		
		\medskip
		\noindent\textbf{Step 1: A single random-walk branch.}
		Fix $p\in\{1,\ldots,P\}$ and write
		\[
		X^*:=x_i+\Delta t\,\mu+\sqrt{\Delta t}\,\sigma\xi^p.
		\]
		
		\smallskip
		\noindent\emph{Case A: $X^*\in\overline\cO$.}
		Here $\widetilde d=0$ and $\widetilde X=X^*$.  Taylor expansion about
		$(t_n,x_i)$ yields
		\begin{align}
			\phi(t_{n+1},\widetilde X)
			={}&\phi_{n,i}+\Delta t\,\partial_t\phi_{n,i}
			+\Delta t\,D\phi_{n,i}\cdot\mu
			+\sqrt{\Delta t}\,D\phi_{n,i}\cdot\sigma\xi^p
			\notag\\
			&+\frac{\Delta t}{2}(\sigma\xi^p)^\top
			D^2\phi_{n,i}(\sigma\xi^p)+O(\Delta t^{3/2}).
			\label{eq:cons-caseA}
		\end{align}
		
		\smallskip
		\noindent\emph{Case B: $X^*\notin\overline\cO$.}
		Let $x_b^p=p^{\gamma_b}(X^*)$, set
		$\widetilde\gamma=\gamma(x_b^p,b)$ and
		$\widetilde d=2d^{\gamma_b}(X^*)$.  Then
		\begin{equation}
			\widetilde X=x_b^p-\frac{\widetilde d}{2}\widetilde\gamma,
			\qquad
			X^*=x_b^p+\frac{\widetilde d}{2}\widetilde\gamma,
			\qquad \widetilde d=O(\sqrt{\Delta t}).
			\label{eq:cons-X-tilde}
		\end{equation}
		Write $\widetilde k=k(t_n,x_b^p,b)$ and
		$\widetilde g=g(t_n,x_b^p,b)$.  Taylor expansion about $(t_n,x_b^p)$ gives
		\begin{align}
			\phi(t_{n+1},\widetilde X)
			={}&\phi(t_n,x_b^p)+\Delta t\,\partial_t\phi(t_n,x_b^p)
			-\frac{\widetilde d}{2}D\phi(t_n,x_b^p)\cdot\widetilde\gamma
			\notag\\
			&+\frac{\widetilde d^2}{8}\widetilde\gamma^\top
			D^2\phi(t_n,x_b^p)\widetilde\gamma+O(\Delta t^{3/2}),
			\label{eq:cons-xb-expand}
		\end{align}
		and
		\begin{equation}
			e^{-\widetilde k\widetilde d}
			=1-\widetilde k\widetilde d+\frac12\widetilde k^2\widetilde d^2
			+O(\widetilde d^3),\qquad
			\widetilde d e^{-\widetilde k\widetilde d/2}
			=\widetilde d-\frac12\widetilde k\widetilde d^2+O(\widetilde d^3).
			\label{eq:cons-exp-expand}
		\end{equation}
		Substitution and collection of powers of $\widetilde d$ yield
		\begin{align}
			&e^{-\widetilde k\widetilde d}\phi(t_{n+1},\widetilde X)
			+\widetilde d e^{-\widetilde k\widetilde d/2}\widetilde g
			\notag\\
			&\quad=\phi(t_n,x_b^p)+\Delta t\,\partial_t\phi(t_n,x_b^p)
			+\widetilde d\left(\widetilde g-\widetilde k\phi(t_n,x_b^p)
			-\frac12D\phi(t_n,x_b^p)\cdot\widetilde\gamma\right)
			\notag\\
			&\qquad+\widetilde d^2\left(
			\frac12\widetilde k^2\phi(t_n,x_b^p)-\frac12\widetilde k\widetilde g
			+\frac12\widetilde kD\phi(t_n,x_b^p)\cdot\widetilde\gamma
			+\frac18\widetilde\gamma^\top D^2\phi(t_n,x_b^p)\widetilde\gamma
			\right)
			+O(\Delta t^{3/2}).
			\label{eq:cons-combined}
		\end{align}
		Adding and subtracting the Taylor terms that move from
		$x_b^p-\widetilde d\widetilde\gamma/2$ to
		$x_b^p+\widetilde d\widetilde\gamma/2$ gives
		\begin{align}
			&e^{-\widetilde k\widetilde d}\phi(t_{n+1},\widetilde X)
			+\widetilde d e^{-\widetilde k\widetilde d/2}\widetilde g
			\notag\\
			&\quad=\Bigl[\phi(t_n,x_b^p)+\Delta t\,\partial_t\phi(t_n,x_b^p)
			+\frac{\widetilde d}{2}D\phi(t_n,x_b^p)\cdot\widetilde\gamma
			+\frac{\widetilde d^2}{8}\widetilde\gamma^\top
			D^2\phi(t_n,x_b^p)\widetilde\gamma\Bigr]
			\notag\\
			&\qquad-\widetilde d\,
			\mathcal L(t_n,x_b^p,D\phi(t_n,x_b^p),\phi(t_n,x_b^p),b)
			\notag\\
			&\qquad+\widetilde d^2\,
			\mathcal R(t_n,x_b^p,D\phi(t_n,x_b^p),\phi(t_n,x_b^p),b)
			+O(\Delta t^{3/2}).
			\label{eq:cons-extract}
		\end{align}
		
		It remains to express the bracketed Taylor polynomial at the starting point.
		From~\eqref{eq:cons-X-tilde},
		\begin{equation}
			x_b^p-x_i=\Delta t\,\mu+\sqrt{\Delta t}\,\sigma\xi^p
			-\frac{\widetilde d}{2}\widetilde\gamma.
			\label{eq:xb-xi-relation}
		\end{equation}
		Set $\delta=x_b^p-x_i$.  Taylor expansion around $x_i$ gives
		\begin{align*}
			\phi(t_n,x_b^p)
			&=\phi_{n,i}+D\phi_{n,i}\cdot\delta
			+\frac12\delta^\top D^2\phi_{n,i}\delta+O(|\delta|^3),\\
			\Delta t\,\partial_t\phi(t_n,x_b^p)
			&=\Delta t\,\partial_t\phi_{n,i}+O(\Delta t|\delta|),\\
			\frac{\widetilde d}{2}D\phi(t_n,x_b^p)\cdot\widetilde\gamma
			&=\frac{\widetilde d}{2}D\phi_{n,i}\cdot\widetilde\gamma
			+\frac{\widetilde d}{2}(D^2\phi_{n,i}\delta)\cdot\widetilde\gamma
			+O(\widetilde d|\delta|^2),\\
			\frac{\widetilde d^2}{8}\widetilde\gamma^\top
			D^2\phi(t_n,x_b^p)\widetilde\gamma
			&=\frac{\widetilde d^2}{8}\widetilde\gamma^\top
			D^2\phi_{n,i}\widetilde\gamma+O(\widetilde d^2|\delta|).
		\end{align*}
		To see the cancellation explicitly, put
		$q=\Delta t\,\mu+\sqrt{\Delta t}\,\sigma\xi^p$, so that
		$\delta=q-\widetilde d\widetilde\gamma/2$.  Then
		\[
		D\phi_{n,i}\cdot\delta
		+\frac{\widetilde d}{2}D\phi_{n,i}\cdot\widetilde\gamma
		=D\phi_{n,i}\cdot q,
		\]
		and, with $H=D^2\phi_{n,i}$,
		\[
		\frac12\delta^\top H\delta
		+\frac{\widetilde d}{2}(H\delta)\cdot\widetilde\gamma
		+\frac{\widetilde d^2}{8}\widetilde\gamma^\top H\widetilde\gamma
		=\frac12q^\top Hq.
		\]
		The drift--drift part of $q^\top Hq$ is $O(\Delta t^2)$ and its mixed
		drift--diffusion part is $O(\Delta t^{3/2})$.  We therefore obtain
		\begin{align}
			&e^{-\widetilde k\widetilde d}\phi(t_{n+1},\widetilde X)
			+\widetilde d e^{-\widetilde k\widetilde d/2}\widetilde g
			\notag\\
			&\quad=\phi_{n,i}+\Delta t\,\partial_t\phi_{n,i}
			+\Delta t\,D\phi_{n,i}\cdot\mu
			+\sqrt{\Delta t}\,D\phi_{n,i}\cdot\sigma\xi^p
			\notag\\
			&\qquad+\frac{\Delta t}{2}(\sigma\xi^p)^\top
			D^2\phi_{n,i}(\sigma\xi^p)
			\notag\\
			&\qquad-\widetilde d_{n,i}^p(a,b)
			\widetilde{\mathcal L}_{n,i}^p
			\bigl(D\phi(t_n,x_b^p),\phi(t_n,x_b^p),a,b\bigr)
			\notag\\
			&\qquad+\bigl(\widetilde d_{n,i}^p(a,b)\bigr)^2
			\widetilde{\mathcal R}_{n,i}^p
			\bigl(D\phi(t_n,x_b^p),\phi(t_n,x_b^p),a,b\bigr)
			+O(\Delta t^{3/2}).
			\label{eq:cons-caseB-final}
		\end{align}
		For a non-exiting branch, the last two lines vanish by definition and
		\eqref{eq:cons-caseB-final} reduces to~\eqref{eq:cons-caseA}.  Thus
		\eqref{eq:cons-caseB-final} is valid for every branch.
		
		\medskip
		\noindent\textbf{Step 2: Average the $P$ equally probable branches.}
		Using
		\[
		\frac1P\sum_{p=1}^P\xi^p=0,
		\qquad
		\frac1P\sum_{p=1}^P\xi^p(\xi^p)^\top=I_m,
		\]
		the linear stochastic term vanishes and the quadratic stochastic term
		becomes a trace.  Therefore
		\begin{align}
			&\frac1P\sum_{p=1}^P\Bigl[
			e^{-\widetilde k_{n,i}^p\widetilde d_{n,i}^p}
			\phi(t_{n+1},\widetilde X_{n,i}^p)
			+\widetilde d_{n,i}^p
			e^{-\widetilde k_{n,i}^p\widetilde d_{n,i}^p/2}
			\widetilde g_{n,i}^p\Bigr]
			\notag\\
			&\quad=\phi_{n,i}+\Delta t\,\partial_t\phi_{n,i}
			+\Delta t\,\mu\cdot D\phi_{n,i}
			+\frac{\Delta t}{2}\Tr(\sigma\sigma^\top D^2\phi_{n,i})
			\notag\\
			&\qquad-\frac1P\sum_{p=1}^P\widetilde d_{n,i}^p
			\widetilde{\mathcal L}_{n,i}^p
			\bigl(D\phi(t_n,x_b^p),\phi(t_n,x_b^p),a,b\bigr)
			\notag\\
			&\qquad+\frac1P\sum_{p=1}^P(\widetilde d_{n,i}^p)^2
			\widetilde{\mathcal R}_{n,i}^p
			\bigl(D\phi(t_n,x_b^p),\phi(t_n,x_b^p),a,b\bigr)
			+O(\Delta t^{3/2}),
			\label{eq:cons-averaged}
		\end{align}
		where the common arguments $(a,b)$ in the branch quantities have only been
		suppressed to shorten the display.
		
		\medskip
		\noindent\textbf{Step 3: Add the ordinary-time discount and the
			$\mathbb P_1$ interpolation error.}
		The interpolation estimate gives, uniformly in the branches,
		\[
		I_h[\phi(t_{n+1},\cdot)|_{\cG_h}](\widetilde X_{n,i}^p)
		=\phi(t_{n+1},\widetilde X_{n,i}^p)+O(h^2),
		\]
		and
		\[
		\frac1{1+r\Delta t}=1-r\Delta t+O(\Delta t^2).
		\]
		Substitution of~\eqref{eq:cons-averaged} into
		\eqref{eq:fixed-operator}, followed by addition of $f\Delta t$ and
		subtraction of $\phi_{n,i}$, gives~\eqref{eq:reflected-expansion-fixed}.
		Indeed, the terms containing $\partial_t\phi$, $\mu$, $\sigma$, $r$, and
		$f$ combine exactly into $-\Delta t[-\partial_t\phi+\mathcal H]$.
		The product of $r\Delta t$ with the boundary sum is
		$O(\Delta t^{3/2})$ because
		$\widetilde d_{n,i}^p=O(\sqrt{\Delta t})$.
		All remainders are uniform on the compact control sets.  Taking the infimum
		over $(a,b)$ and using
		\(
		|\inf F_{a,b}-\inf G_{a,b}|\le\sup|F_{a,b}-G_{a,b}|
		\)
		and $\inf(-F)=-\sup F$ proves
		\eqref{eq:reflected-expansion-controlled}.
		
		\medskip
		\noindent\textbf{Step 4: Boundary nodes.}
		Fix $i\in\mathcal I_h^\partial$ and $b\in B$, and set
		\[
		y_{i,b}=x_i-\ell_h\gamma(x_i,b),\qquad
		\phi_n=\phi(t_n,\cdot)|_{\cG_h}.
		\]
		The $\mathbb P_1$ interpolation estimate and Taylor's formula at $x_i$
		give
		\[
		I_h[\phi_n](y_{i,b})
		=\phi_{n,i}-\ell_h\gamma(x_i,b)\cdot D\phi_{n,i}
		+O(\ell_h^2+h^2).
		\]
		Substitution into~\eqref{eq:fixed-boundary-map} yields, uniformly in $b$,
		\begin{align}
			\cS_{n,i}^{\partial,b}[\phi_n]-\phi_{n,i}
			&=-\frac{\ell_h}{1+\ell_h k(t_n,x_i,b)}
			\mathcal L(t_n,x_i,D\phi_{n,i},\phi_{n,i},b)
			+O(\ell_h^2+h^2)
			\notag\\
			&=-\ell_h\mathcal L(t_n,x_i,D\phi_{n,i},\phi_{n,i},b)+O(h^2),
			\label{eq:boundary-fixed-consistency}
		\end{align}
		where the second equality uses $\ell_h\asymp h$ and boundedness of $k$ and
		$\mathcal L$.  Taking the infimum over $b$ and using uniformity of the
		remainder proves~\eqref{eq:boundary-controlled-consistency}.
	\end{proof}
	
	\subsection{Stability of the parabolic scheme}
	
	\begin{proposition}[Parabolic stability]\label{prop:stability}
		Under Assumptions~\ref{ass:H1}--\ref{ass:H4}, assume
		\begin{equation}
			\Delta t,h\to0,
			\qquad \frac{h^2}{\Delta t}\to0.
			\label{eq:admissible-refinement}
		\end{equation}
		There is a constant $C$, independent of the discretization parameters, such
		that
		\[
		\max_{0\le n\le N_T}\|U^n\|_\infty\le C
		\]
		for all sufficiently fine grids.
	\end{proposition}
	\begin{proof}
		Set
		\[
		F_0=\|f\|_\infty,\qquad G_0=\|g\|_\infty,\qquad A_0=G_0+1.
		\]
		
		\medskip
		\noindent\textbf{Step 1: Auxiliary function.}
		By uniform obliqueness,
		\[
		\nu:=\min_{(x,b)\in\partial\cO\times B}n(x)\cdot\gamma(x,b)>0.
		\]
		Let $d_\partial(x)=\dist(x,\partial\cO)$ in $\cO$ and write
		$L_\rho:=\{x\in\overline\cO:d_\partial(x)<\rho\}$.  Choose a tubular
		radius $\eta>0$ on which $d_\partial$ is $C^3$ and
		$Dd_\partial=-n$ on $\partial\cO$, put $\kappa=1/\nu$, and choose
		$\varepsilon_0\le\min\{\eta/2,1/(2\kappa)\}$.  On $L_{2\varepsilon_0}$ set
		$\chi_0=1-\kappa d_\partial$.  Let $\zeta$ be a smooth cutoff equal to one
		on $L_{\varepsilon_0}$ and zero outside $L_{2\varepsilon_0}$, and extend
		$\widetilde\chi:=\zeta\chi_0$ by zero to the rest of $\overline\cO$.
		Set $\chi:=\widetilde\chi+\|\widetilde\chi^-\|_\infty$.  Then
		\begin{equation}
			\chi\in C^3(\overline\cO),\qquad \chi\ge0,\qquad
			\gamma(x,b)\cdot D\chi(x)\ge1
			\quad (x,b)\in\partial\cO\times B.
			\label{eq:chi}
		\end{equation}
		The cutoff can be chosen flat where it vanishes; hence $\chi$ has a bounded
		$C^3$ extension to a fixed neighbourhood of $\overline\cO$.
		
		\medskip
		\noindent\textbf{Step 2: Upper barrier.}
		For a constant $M_0>0$ to be fixed below, define
		\[
		\psi(t,x)=e^{T-t}(M_0+A_0\chi(x)).
		\]
		Set
		\[
		C_1:=\sup_{(t,x,a)\in[0,T]\times\overline\cO\times A}
		\left|\frac12\Tr(\sigma\sigma^\top D^2\chi)
		+\mu\cdot D\chi\right|<\infty.
		\]
		Since $r,k\ge0$ and $e^{T-t}\ge1$, direct substitution gives, uniformly
		in $a\in A$ and $b\in B$,
		\begin{align}
			[-\partial_t\psi+\mathcal H]
			(t,x,D\psi,D^2\psi,\psi,a)
			&\ge e^{T-t}M_0-e^TA_0C_1-F_0,
			\label{eq:H-psi}\\
			\mathcal L(t,x,D\psi,\psi,b)&\ge A_0-G_0=1
			\qquad (x\in\partial\cO),
			\label{eq:L-psi}
		\end{align}
		Choose
		\begin{equation}
			M_0\ge\max\{e^TA_0C_1+F_0+4,\ \|\Psi\|_\infty\}.
			\label{eq:M0}
		\end{equation}
		Then the two residuals in \eqref{eq:H-psi}--\eqref{eq:L-psi} are bounded
		below by $4$ and $1$, respectively.
		
		\medskip
		\noindent\textbf{Step 3: Consistency estimate for the barriers.}
		Apply Proposition~\ref{prop:nodewise-consistency} to the two fixed smooth
		functions $\psi$ and $-\psi$, and let $C_{\rm cons}$ be a common constant
		for their interior and boundary remainders.  By
		\eqref{eq:admissible-refinement} and $\ell_h\ge c_0h$, for all sufficiently
		fine grids,
		\begin{equation}
			C_{\rm cons}(\Delta t^{3/2}+h^2)\le\Delta t,
			\qquad C_{\rm cons}h^2\le\frac12\ell_h.
			\label{eq:barrier-remainder}
		\end{equation}
		
		On an exiting branch,
		\eqref{eq:R-equals-kL} gives
		\[
		-\widetilde d\,\widetilde{\mathcal L}
		+\widetilde d^2\widetilde{\mathcal R}
		=-\widetilde d\left(1-\frac12\widetilde k\widetilde d\right)
		\widetilde{\mathcal L}.
		\]
		The step bound~\eqref{eq:step-bound} implies
		$\widetilde d=O(\sqrt{\Delta t})$ uniformly.  Hence, for sufficiently small
		$\Delta t$,
		\(
		1-\widetilde k\widetilde d/2\ge1/2
		\)
		on every branch.  Part~\textup{(I)} of
		Proposition~\ref{prop:nodewise-consistency} can therefore be used with a
		fixed sign for every reflected boundary contribution.
		
		\medskip
		\noindent\textbf{Step 4: The upper barrier is a discrete supersolution.}
		Fix $i\in\mathcal I_h^\circ$ and $(a,b)\in A\times B$.  Part~\textup{(I)}
		of Proposition~\ref{prop:nodewise-consistency} gives
		\begin{align}
			&\cS_{n,i}^{a,b}[\psi(t_{n+1},\cdot)]-\psi(t_n,x_i)
			\notag\\
			&\quad=-\Delta t\,
			[-\partial_t\psi+\mathcal H](t_n,x_i,D\psi,D^2\psi,\psi,a)
			\notag\\
			&\qquad\quad+\frac1P\sum_{p=1}^P\left\{
			-\widetilde d_{n,i}^p\,
			\widetilde{\mathcal L}_{n,i}^p[\psi](a,b)
			+(\widetilde d_{n,i}^p)^2
			\widetilde{\mathcal R}_{n,i}^p[\psi](a,b)\right\}
			+\rho_{n,i}^{a,b},
			\label{eq:upper-fixed-expansion}
		\end{align}
		Here $\widetilde{\mathcal L}_{n,i}^p[\psi](a,b)$ and
		$\widetilde{\mathcal R}_{n,i}^p[\psi](a,b)$ denote the two truncated
		operators in Proposition~\ref{prop:nodewise-consistency} evaluated at
		$\psi$; they are zero on a non-exiting branch.  Moreover,
		$|\rho_{n,i}^{a,b}|\le C_{\rm cons}(\Delta t^{3/2}+h^2)\le\Delta t$.
		The first term on the right of~\eqref{eq:upper-fixed-expansion} is at most
		$-4\Delta t$.  A non-exiting branch has $\widetilde d_{n,i}^p=0$.
		On an exiting branch, \eqref{eq:R-equals-kL}, \eqref{eq:L-psi}, and
		$1-\widetilde k_{n,i}^p\widetilde d_{n,i}^p/2\ge1/2$ give
		\[
		-\widetilde d_{n,i}^p\widetilde{\mathcal L}_{n,i}^p[\psi]
		+(\widetilde d_{n,i}^p)^2\widetilde{\mathcal R}_{n,i}^p[\psi]
		=-\widetilde d_{n,i}^p
		\left(1-\frac12\widetilde k_{n,i}^p\widetilde d_{n,i}^p\right)
		\widetilde{\mathcal L}_{n,i}^p[\psi]
		\le-\frac12\widetilde d_{n,i}^p\le0.
		\]
		Consequently, for every fixed $(a,b)$,
		\[
		\cS_{n,i}^{a,b}[\psi(t_{n+1},\cdot)]-\psi(t_n,x_i)
		\le-4\Delta t-\frac1{2P}\sum_{p:\,\mathrm{exit}}
		\widetilde d_{n,i}^p+\Delta t
		\le-3\Delta t.
		\]
		Taking the infimum over $(a,b)$ proves
		\begin{equation}
			S_{n,i}[\psi(t_{n+1},\cdot)]-\psi(t_n,x_i)\le-3\Delta t<0.
			\label{eq:S-psi-interior}
		\end{equation}
		
		Now fix $i\in\mathcal I_h^\partial$ and $b\in B$.  The fixed-control
		boundary expansion~\eqref{eq:boundary-fixed-consistency} gives
		\[
		\cS_{n,i}^{\partial,b}[\psi(t_n,\cdot)]-\psi(t_n,x_i)
		=-\ell_h\mathcal L(t_n,x_i,D\psi,\psi,b)
		+\rho_{n,i}^{\partial,b},
		\qquad |\rho_{n,i}^{\partial,b}|\le C_{\rm cons}h^2.
		\]
		Since \eqref{eq:L-psi} holds for every $b$,
		\[
		\cS_{n,i}^{\partial,b}[\psi(t_n,\cdot)]-\psi(t_n,x_i)
		\le-\ell_h+C_{\rm cons}h^2\le-\frac12\ell_h.
		\]
		Taking the infimum over $b$ proves
		\begin{equation}
			S_{n,i}^{\partial}[\psi(t_n,\cdot)]-\psi(t_n,x_i)
			\le-\frac12\ell_h<0.
			\label{eq:S-psi-boundary}
		\end{equation}
		Thus \eqref{eq:S-psi-interior} and \eqref{eq:S-psi-boundary} verify the
		upper barrier separately on the interior and boundary node sets.  Notice
		that the boundary estimate uses only $\ell_h\asymp h$ and no time-step
		condition.
		
		\medskip
		\noindent\textbf{Step 5: Lower barrier.}
		Define
		\[
		\underline\psi(t,x):=-e^{T-t}(M_0+A_0\chi(x))=-\psi(t,x).
		\]
		Direct substitution, exactly as in Step~2, gives
		\begin{align}
			[-\partial_t\underline\psi+\mathcal H]
			(t,x,D\underline\psi,D^2\underline\psi,\underline\psi,a)&\le-4,
			\label{eq:H-lower-psi}\\
			\mathcal L(t,x,D\underline\psi,\underline\psi,b)&\le-1
			\qquad(x\in\partial\cO).
			\label{eq:L-lower-psi}
		\end{align}
		Fix an interior node and $(a,b)$.  Substitution of $\underline\psi$ into
		the fixed-control expansion~\eqref{eq:upper-fixed-expansion}, with a new
		remainder satisfying the same absolute bound, gives a volume contribution
		at least $4\Delta t$.  On every exiting branch,
		\[
		-\widetilde d\left(1-\frac12\widetilde k\widetilde d\right)
		\widetilde{\mathcal L}[\underline\psi]
		\ge\frac12\widetilde d,
		\]
		because $\widetilde{\mathcal L}[\underline\psi]\le-1$.  Hence, for every
		fixed $(a,b)$,
		\[
		\cS_{n,i}^{a,b}[\underline\psi(t_{n+1},\cdot)]
		-\underline\psi(t_n,x_i)
		\ge4\Delta t+\frac1{2P}\sum_{p:\,\mathrm{exit}}
		\widetilde d_{n,i}^p-\Delta t
		\ge3\Delta t.
		\]
		Because this lower bound holds for every control pair, taking the infimum
		preserves it and yields
		\begin{equation}
			S_{n,i}[\underline\psi(t_{n+1},\cdot)]
			-\underline\psi(t_n,x_i)\ge3\Delta t>0,
			\qquad i\in\mathcal I_h^\circ.
			\label{eq:S-lower-psi-interior}
		\end{equation}
		At a boundary node, the fixed-control boundary expansion and
		\eqref{eq:L-lower-psi} give, for every $b$,
		\[
		\cS_{n,i}^{\partial,b}[\underline\psi(t_n,\cdot)]
		-\underline\psi(t_n,x_i)
		\ge\ell_h-C_{\rm cons}h^2\ge\frac12\ell_h.
		\]
		Taking the infimum over $b$ therefore gives
		\begin{equation}
			S_{n,i}^{\partial}[\underline\psi(t_n,\cdot)]
			-\underline\psi(t_n,x_i)
			\ge\frac12\ell_h>0.
			\label{eq:S-lower-psi-boundary}
		\end{equation}
		
		\medskip
		\noindent\textbf{Step 6: Comparison and the $L^\infty$ estimate.}
		At $t=T$, the choice of $M_0$ gives
		\[
		\underline\psi(T,x_i)\le U_i^{N_T}\le\psi(T,x_i)
		\qquad(i\in\mathcal I_h).
		\]
		Assume this order at level $n+1$.  Interior monotonicity together with
		\eqref{eq:S-psi-interior}--\eqref{eq:S-lower-psi-interior} first yields
		\[
		\underline\psi(t_n,x_i)\le U_i^n\le\psi(t_n,x_i)
		\qquad(i\in\mathcal I_h^\circ).
		\]
		For the same time layer, write $T_n^q$ for the boundary map with interior
		vector $q$.  It is order preserving in both arguments and is a contraction.
		The just-proved interior bounds and
		\eqref{eq:S-psi-boundary}--\eqref{eq:S-lower-psi-boundary} give
		\[
		T_n^{U_\circ^n}(\psi_\partial^n)
		\le T_n^{\psi_\circ^n}(\psi_\partial^n)\le\psi_\partial^n,
		\qquad
		T_n^{U_\circ^n}(\underline\psi_\partial^n)
		\ge T_n^{\underline\psi_\circ^n}(\underline\psi_\partial^n)
		\ge\underline\psi_\partial^n.
		\]
		Iterating the monotone contraction $T_n^{U_\circ^n}$ from the two barriers
		therefore brackets its unique fixed point between them.  Hence the same bound
		holds on $\mathcal I_h^\partial$.  Backward induction proves
		\[
		\max_{0\le n\le N_T}\|U^n\|_\infty
		\le e^T(M_0+A_0\|\chi\|_\infty),
		\]
		which is independent of $h$ and $\Delta t$.
	\end{proof}
	\begin{remark}\upshape
		Condition~\eqref{eq:admissible-refinement} is used in the barrier proof to
		absorb the $O(h^2)$ consistency remainder into the $O(\Delta t)$ volume
		margin.  It is a refinement condition for the present stability and
		convergence results, not a CFL restriction for positivity or monotonicity;
		we do not claim the stability bound for arbitrary independent choices of
		$h$ and $\Delta t$.
	\end{remark}
	\section{Convergence to the viscosity solution}
	\label{sec:convergence}
	
	For $t\in[t_n,t_{n+1})$ set
	\begin{equation}
		u^{h,\dt}(t,x)=I_h[U^n](x),
		\qquad
		u^{h,\dt}(T,x)=I_h[U^{N_T}](x).
		\label{eq:extension}
	\end{equation}
	All limits in this section are taken along families satisfying~\eqref{eq:admissible-refinement}.  Define the nodal half-relaxed limits
	\begin{align}
		\overline u(t,x)
		&=\limsup_{\substack{h,\dt\to0,\ h^2/\dt\to0\\
				(t_n,x_i)\to(t,x)}}U_i^n,
		\label{eq:upper-relaxed}\\
		\underline u(t,x)
		&=\liminf_{\substack{h,\dt\to0,\ h^2/\dt\to0\\
				(t_n,x_i)\to(t,x)}}U_i^n.
		\label{eq:lower-relaxed}
	\end{align}
	They are also the half-relaxed limits of~\eqref{eq:extension}: every interpolated value is a convex combination of nodal values at distance $O(h)$, whereas every nodal value is itself an interpolated value.  Proposition~\ref{prop:stability} makes both limits finite.
	
	\begin{theorem}[Uniform viscosity convergence]\label{thm:convergence}
		Under Assumptions~\ref{ass:H1}--\ref{ass:H4}, let~\eqref{eq:admissible-refinement} hold.  Then
		$u^{h,\dt}$ converges uniformly on $\overline\cO_T$ to the unique viscosity solution of~\eqref{eq:HJB-main}.
	\end{theorem}
	
	\begin{proof}
		By the smooth-test reduction~\cite[Sections~2--3]{CrandallIshiiLions1992},
		it is enough to use test functions having a bounded $C^3$ extension to a
		fixed neighbourhood of the contact point.  Multiplication by a cutoff equal
		to one near that point does not change the argument: the interior stencil has
		diameter $O(\sqrt{\dt}+h)$ by~\eqref{eq:step-bound}, and the same-level
		boundary stencil has diameter $O(h)$.
		
		\medskip
		\noindent\textbf{Step 1: Strict maximum and recovery sequence.}
		Let $z_0=(t_0,x_0)\in\overline\cO_T$, and suppose that
		$\overline u-\phi$ has a local maximum at $z_0$.  Subtracting a constant,
		assume $(\overline u-\phi)(z_0)=0$.  For fixed $\varepsilon>0$, set
		\begin{equation}
			\phi^\varepsilon(t,x)
			:=\phi(t,x)+\varepsilon\bigl(|t-t_0|^2+|x-x_0|^4\bigr).
			\label{eq:strict-upper-test}
		\end{equation}
		Then $\overline u-\phi^\varepsilon$ has a strict local maximum $0$ at
		$z_0$, and the derivatives of $\phi^\varepsilon$ used in the equation agree
		with those of $\phi$ at $z_0$.
		
		Choose compact relative neighbourhoods
		$K'\Subset K\subset\overline\cO_T$ such that the strict maximum lies in
		$K'$ and $\overline u-\phi^\varepsilon<0$ on $K\setminus K'$.  On a
		subsequence of discretizations realizing the upper relaxed limit, let
		$\mathcal G_j^{\rm st}$ be the corresponding space--time nodal grid and
		write $U_j(t_n,x_i):=U_i^n$.  Set
		\begin{equation}
			c_j:=\max_{(t_n,x_i)\in K\cap\mathcal G_j^{\rm st}}
			\{U_i^n-\phi^\varepsilon(t_n,x_i)\},
			\qquad
			z_j=(t_{n_j},x_{i_j})
			\label{eq:upper-grid-maximum}
		\end{equation}
		at a maximizing node.  A recovery sequence
		$y_j\in K\cap\mathcal G_j^{\rm st}$ with
		$y_j\to z_0$ and $U_j(y_j)\to\overline u(z_0)$ gives
		\[
		\liminf_{j\to\infty}c_j
		\ge\lim_{j\to\infty}\bigl(U_j(y_j)-\phi^\varepsilon(y_j)\bigr)=0.
		\]
		By stability, every subsequence of $(z_j,c_j)$ has a further subsequence
		with $z_j\to z_*\in K$ and $c_j\to c\ge0$.  Since
		$U_{i_j}^{n_j}=\phi^\varepsilon(z_j)+c_j$, the definition of
		$\overline u$ gives
		\[
		(\overline u-\phi^\varepsilon)(z_*)\ge c\ge0.
		\]
		Strictness forces $z_*=z_0$ and $c=0$.  This argument applies to every
		subsequence; hence
		\begin{equation}
			z_j\to z_0,\qquad c_j\to0.
			\label{eq:contact-limit}
		\end{equation}
		For large $j$, $z_j\in K'$ and every stencil at $z_j$ is contained in
		$K$.  Therefore the maximizing property yields
		\begin{equation}
			U^{n_j+1}\le
			\phi^\varepsilon(t_{n_j+1},\cdot)|_{\cG_{h_j}}+c_j
			\quad\hbox{on every next-time stencil},
			\label{eq:max-next-stencil}
		\end{equation}
		whenever $n_j<N_{T,j}$, and
		\begin{equation}
			U^{n_j}\le
			\phi^\varepsilon(t_{n_j},\cdot)|_{\cG_{h_j}}+c_j
			\quad\hbox{on every same-time boundary stencil}.
			\label{eq:max-boundary-stencil}
		\end{equation}
		
		\medskip
		\noindent\textbf{Step 2: Monotonicity and signed constant shifts.}
		The constants $c_j$ need not have a fixed sign.  For an interior
		fixed-control operator, define
		\[
		\delta_{j}^{a,b}(c):=
		\cS_{n_j,i_j}^{a,b}[V+c]-\cS_{n_j,i_j}^{a,b}[V]-c.
		\]
		By~\eqref{eq:shift-defect}, uniformly in the controls,
		\begin{equation}
			|\delta_j^{a,b}(c)|
			\le C|c|\left(\dt_j+\frac1P\sum_{p=1}^P
			\widetilde d_{n_j,i_j}^p(a,b)\right).
			\label{eq:interior-signed-shift}
		\end{equation}
		For a boundary fixed-control row the formula is exact:
		\begin{equation}
			\cS_{n_j,i_j}^{\partial,b}[V+c]
			-\cS_{n_j,i_j}^{\partial,b}[V]-c
			=-\frac{\ell_{h_j}k(t_{n_j},x_{i_j},b)}
			{1+\ell_{h_j}k(t_{n_j},x_{i_j},b)}\,c.
			\label{eq:boundary-signed-shift}
		\end{equation}
		Thus, after division by the natural interior scale
		$\dt_j+P^{-1}\sum_p\widetilde d^p$ or by the boundary scale $\ell_{h_j}$,
		the shift errors produced by $c_j$ are $O(|c_j|)$.
		
		If $i_j\in\mathcal I_{h_j}^\circ$ and $n_j<N_{T,j}$, the scheme,
		\eqref{eq:max-next-stencil}, and monotonicity imply, for every fixed
		$(a,b)$,
		\begin{equation}
			0\le
			\cS_{n_j,i_j}^{a,b}[\phi^\varepsilon(t_{n_j+1},\cdot)]
			-\phi^\varepsilon(z_j)+\delta_j^{a,b}(c_j).
			\label{eq:upper-interior-discrete}
		\end{equation}
		If $i_j\in\mathcal I_{h_j}^\partial$, then for every fixed $b$,
		\begin{equation}
			0\le
			\cS_{n_j,i_j}^{\partial,b}[\phi^\varepsilon(t_{n_j},\cdot)]
			-\phi^\varepsilon(z_j)
			-\frac{\ell_{h_j}k(t_{n_j},x_{i_j},b)}
			{1+\ell_{h_j}k(t_{n_j},x_{i_j},b)}c_j.
			\label{eq:upper-boundary-discrete}
		\end{equation}
		Indeed, at an interior node
		$U_{i_j}^{n_j}=S_{n_j,i_j}[U^{n_j+1}]$, whereas at a boundary node
		$U_{i_j}^{n_j}=S_{n_j,i_j}^{\partial}[U^{n_j}]$; these are precisely the
		two lines of~\eqref{eq:scheme}.
		
		\medskip
		\noindent\textbf{Step 3: Interior subsolution condition.}
		Assume $t_0<T$ and $x_0\in\cO$.  For all sufficiently large $j$,
		$i_j\in\mathcal I_{h_j}^\circ$ and every branch stays in $\cO$, uniformly
		in $(a,b)$.  Fix arbitrary $(a,b)\in A\times B$.  In
		\eqref{eq:upper-interior-discrete} all $\widetilde d^p$ vanish.  Part
		\textup{(I)} of Proposition~\ref{prop:nodewise-consistency}, divided by
		$\dt_j$, therefore gives
		\[
		[-\partial_t\phi^\varepsilon+\mathcal H]
		(t_{n_j},x_{i_j},D\phi^\varepsilon,D^2\phi^\varepsilon,
		\phi^\varepsilon,a)
		\le C\left(\sqrt{\dt_j}+\frac{h_j^2}{\dt_j}+|c_j|\right).
		\]
		The right-hand side tends to zero by
		\eqref{eq:admissible-refinement} and~\eqref{eq:contact-limit}.  Passing to
		the limit and using the arbitrariness of $a$ gives
		\[
		-\partial_t\phi^\varepsilon(z_0)
		+H(t_0,x_0,D\phi^\varepsilon(z_0),D^2\phi^\varepsilon(z_0),
		\overline u(z_0))\le0.
		\]
		Letting $\varepsilon\downarrow0$ proves~\eqref{eq:sub-interior}.
		
		\medskip
		\noindent\textbf{Step 4: Boundary subsolution condition.}
		Assume $t_0<T$ and $x_0\in\partial\cO$.  Passing to a subsequence, the
		contact nodes are either all boundary nodes or all interior nodes.
		
		\smallskip
		\noindent\textit{Case A: boundary contact nodes.}
		Let $i_j\in\mathcal I_{h_j}^\partial$ and fix $b\in B$.  Divide
		\eqref{eq:upper-boundary-discrete} by $\ell_{h_j}$.  Part~\textup{(II)} of
		Proposition~\ref{prop:nodewise-consistency},
		\eqref{eq:boundary-signed-shift}, $c_j\to0$, and
		$h_j^2/\ell_{h_j}=O(h_j)$ yield
		\[
		\mathcal L(t_0,x_0,D\phi^\varepsilon(z_0),
		\overline u(z_0),b)\le0.
		\]
		This holds for every $b$, and hence
		\begin{equation}
			L(t_0,x_0,D\phi^\varepsilon(z_0),\overline u(z_0))\le0.
			\label{eq:boundary-node-sub}
		\end{equation}
		
		\smallskip
		\noindent\textit{Case B: interior contact nodes.}
		Let $i_j\in\mathcal I_{h_j}^\circ$, fix $(a,b)\in A\times B$, and abbreviate
		\[
		\begin{aligned}
			d_{j,p}&:=\widetilde d_{n_j,i_j}^p(a,b),&
			k_{j,p}&:=\widetilde k_{n_j,i_j}^p(a,b),&
			x_{j,p}^b&:=x_{n_j,i_j}^{b,p}(a),\\
			s_j&:=\dt_j+\frac1P\sum_{p=1}^P d_{j,p},&
			\lambda_j&:=\frac{\dt_j}{s_j}.&&
		\end{aligned}
		\]
		For an exiting branch set
		\[
		L_{j,p}^b:=\mathcal L(t_{n_j},x_{j,p}^b,
		D\phi^\varepsilon(t_{n_j},x_{j,p}^b),
		\phi^\varepsilon(t_{n_j},x_{j,p}^b),b).
		\]
		The value assigned to $L_{j,p}^b$ on a non-exiting branch is irrelevant
		because then $d_{j,p}=0$.  Using
		$\mathcal R=(k/2)\mathcal L$ in Part~\textup{(I)} of Proposition
		~\ref{prop:nodewise-consistency}, inequality
		\eqref{eq:upper-interior-discrete} becomes
		\begin{equation}
			\lambda_j A_j^a+\frac1{Ps_j}\sum_{p=1}^P d_{j,p}
			\left(1-\frac12k_{j,p}d_{j,p}\right)L_{j,p}^b
			\le\eta_j,
			\label{eq:sub-normalized}
		\end{equation}
		where
		\begin{align*}
			A_j^a&:=[-\partial_t\phi^\varepsilon+\mathcal H]
			(t_{n_j},x_{i_j},D\phi^\varepsilon(z_j),D^2\phi^\varepsilon(z_j),
			\phi^\varepsilon(z_j),a),\\
			|\eta_j|&\le C\left(
			\sqrt{\dt_j}+\frac{h_j^2}{\dt_j}+|c_j|\right)\longrightarrow0.
		\end{align*}
		Here $s_j\ge\dt_j$ controls the consistency remainder, and
		\eqref{eq:interior-signed-shift} controls the last term by $C|c_j|$.
		
		This includes the case in which no branch exits: then every $d_{j,p}=0$,
		$s_j=\dt_j$, and $\lambda_j=1$.  Extract a subsequence such that
		$\lambda_j\to\lambda\in[0,1]$.  Since
		$\max_p d_{j,p}=O(\sqrt{\dt_j})$,
		\begin{equation}
			\frac1{Ps_j}\sum_{p=1}^P d_{j,p}
			\left(1-\frac12k_{j,p}d_{j,p}\right)
			=1-\lambda_j+o(1).
			\label{eq:boundary-weight-limit}
		\end{equation}
		Every exiting projection $x_{j,p}^b$ tends to $x_0$.  Since $k$ is bounded
		and $\max_p d_{j,p}\to0$, the factors
		$1-k_{j,p}d_{j,p}/2$ are nonnegative for all large $j$; hence the absolute
		sum of the weights in~\eqref{eq:boundary-weight-limit} is uniformly
		bounded.  Uniform continuity of $\mathcal L$ therefore gives
		\[
		\frac1{Ps_j}\sum_{p=1}^P d_{j,p}
		\left(1-\frac12k_{j,p}d_{j,p}\right)
		\bigl[L_{j,p}^b-\mathcal L(t_0,x_0,D\phi^\varepsilon(z_0),
		\overline u(z_0),b)\bigr]\longrightarrow0.
		\]
		Passing to the limit in~\eqref{eq:sub-normalized} yields
		\begin{equation}
			\lambda[-\partial_t\phi^\varepsilon+\mathcal H]
			(t_0,x_0,D\phi^\varepsilon(z_0),D^2\phi^\varepsilon(z_0),
			\overline u(z_0),a)
			+(1-\lambda)\mathcal L(t_0,x_0,D\phi^\varepsilon(z_0),
			\overline u(z_0),b)\le0.
			\label{eq:boundary-convex-sub}
		\end{equation}
		If both $-\partial_t\phi^\varepsilon+H$ and $L$ were positive at $z_0$,
		compactness of $A$ and $B$ and continuity of the control-resolved operators
		would give fixed controls $a_0\in A$, $b_0\in B$, and $\delta>0$ such that
		both corresponding quantities are at least $\delta$.  Repeating the
		preceding subsequence extraction for $(a_0,b_0)$, the left-hand side of
		\eqref{eq:boundary-convex-sub} would be at least
		$\lambda\delta+(1-\lambda)\delta=\delta>0$, a contradiction.  Thus
		\[
		\min\{-\partial_t\phi^\varepsilon+H,L\}(z_0)\le0.
		\]
		Together with Case~A, and then letting $\varepsilon\downarrow0$, this proves
		\eqref{eq:sub-boundary}.  If boundary and interior contact nodes alternate,
		one takes an infinite subsequence of either type; there is no third case.
		
		\medskip
		\noindent\textbf{Step 5: Terminal subsolution condition.}
		Let $t_0=T$.  Repeat Step~1 with the strict test
		$\phi^\varepsilon$ and obtain $z_j=(t_{n_j},x_{i_j})\to(T,x_0)$ and
		$c_j\to0$.
		
		If $t_{n_j}=T$ along an infinite subsequence, then
		\[
		U_{i_j}^{N_{T,j}}=\Psi(x_{i_j}),\qquad
		U_{i_j}^{N_{T,j}}=\phi^\varepsilon(T,x_{i_j})+c_j.
		\]
		Passing to the limit gives
		$\overline u(T,x_0)=\phi(T,x_0)=\Psi(x_0)$, so the terminal-data member of
		the relaxed minimum is zero.
		
		Otherwise $t_{n_j}<T$ for all sufficiently large $j$.  The next time level
		may be $T$, but Steps~2--4 remain valid.  Letting $\varepsilon\downarrow0$,
		if $x_0\in\cO$ they give
		$-\phi_t+H\le0$; if $x_0\in\partial\cO$ they give
		$\min\{-\phi_t+H,L\}\le0$.  Consequently,
		\[
		\min\{-\phi_t+H,\overline u-\Psi\}(T,x_0)\le0
		\quad(x_0\in\cO),
		\]
		and
		\[
		\min\{-\phi_t+H,L,\overline u-\Psi\}(T,x_0)\le0
		\quad(x_0\in\partial\cO).
		\]
		Thus $\overline u$ is a viscosity subsolution.
		
		\medskip
		\noindent\textbf{Step 6: Supersolution inequalities.}
		Let $\underline u-\phi$ have a local minimum $0$ at
		$z_0=(t_0,x_0)$ and set
		\[
		\phi_\varepsilon(t,x)
		:=\phi(t,x)-\varepsilon\bigl(|t-t_0|^2+|x-x_0|^4\bigr).
		\]
		For brevity, denote the control-resolved residuals at $z_0$ by
		\begin{align*}
			A_\varepsilon^a
			&:=[-\partial_t\phi_\varepsilon+\mathcal H]
			(t_0,x_0,D\phi_\varepsilon(z_0),D^2\phi_\varepsilon(z_0),
			\underline u(z_0),a),\\
			L_\varepsilon^b
			&:=\mathcal L(t_0,x_0,D\phi_\varepsilon(z_0),
			\underline u(z_0),b).
		\end{align*}
		Repeating Step~1 with minima in place of maxima gives recovery nodes
		$z_j\to z_0$, constants $c_j\to0$, and the corresponding lower stencil
		comparisons.  Since both discrete updates are infima, these comparisons are
		applied below to controls chosen within $\dt_j^2$ (respectively
		$\ell_{h_j}^2$) of the relevant infimum.
		For $t_0<T$ and $x_0\in\cO$, choose controls within $\dt_j^2$ of the
		interior infimum and repeat Step~3.  After division by $\dt_j$, the remainder
		is $O(\sqrt{\dt_j}+h_j^2/\dt_j+|c_j|+\dt_j)=o(1)$, and compactness gives a
		control $a_*$ for which $A_\varepsilon^{a_*}\ge0$; hence
		$-\partial_t\phi_\varepsilon+H\ge0$.
		
		Let now $t_0<T$ and $x_0\in\partial\cO$.  A subsequence of boundary
		contact nodes is treated by Case~A with an $\ell_{h_j}^2$-minimizer of the
		boundary infimum, and gives $L_\varepsilon^{b_*}\ge0$, hence $L\ge0$.
		For interior contact nodes approaching the boundary, choose $(a_j,b_j)$
		so that the fixed-control value is within $\dt_j^2$ of the interior
		infimum, and use the same
		$s_j=\dt_j+P^{-1}\sum_p d_{j,p}$ and $\lambda_j=\dt_j/s_j$ as in Case~B.
		Since $s_j\ge\dt_j$, the normalized selection error satisfies
		$\dt_j^2/s_j\le\dt_j\to0$.  After extraction,
		$(a_j,b_j)\to(a_*,b_*)$ and $\lambda_j\to\lambda\in[0,1]$.
		The estimates in Case~B are uniform in the controls; with $d_{j,p}=0$ on
		non-exiting branches, they also cover changes in the exiting-branch set.
		The reversed Case~B inequality and~\eqref{eq:boundary-weight-limit} then give
		\[
		\lambda A_\varepsilon^{a_*}+(1-\lambda)L_\varepsilon^{b_*}\ge0,
		\qquad 0\le\lambda\le1.
		\]
		Thus at least one residual is nonnegative and
		$\max\{-\partial_t\phi_\varepsilon+H,L\}(z_0)\ge0$.  Alternating contact
		types are covered by passing to an infinite subsequence.  This implication
		also includes the endpoint cases $\lambda=0$ and $\lambda=1$.
		
		At $t_0=T$, the same construction either reaches the terminal layer and
		gives $\underline u-\Psi\ge0$, or reduces to the preceding nonterminal
		argument.  Letting $\varepsilon\downarrow0$ yields the terminal relaxed
		maximum inequalities in Definition~\ref{def:viscosity}.  Therefore
		$\underline u$ is a viscosity supersolution.
		
		\medskip
		\noindent\textbf{Step 7: Comparison and uniform convergence.}
		By definition, $\underline u\le\overline u$.  Steps~1--6 and
		Proposition~\ref{thm:wellposed} give $\overline u\le\underline u$.
		Hence $u:=\overline u=\underline u$ is continuous and is a viscosity
		solution of~\eqref{eq:HJB-main}; the same comparison principle gives
		uniqueness.  The remainder of the argument proves uniform convergence to
		this $u$.
		
		If convergence were not uniform, compactness of $\overline\cO_T$ would
		give $\varepsilon_0>0$, a subsequence, and points $z_j\to z_*$ with
		$|u^{h_j,\dt_j}(z_j)-u(z_j)|\ge\varepsilon_0$.  A positive discrepancy
		would imply $\overline u(z_*)\ge u(z_*)+\varepsilon_0$; a negative one would
		imply $\underline u(z_*)\le u(z_*)-\varepsilon_0$.  Both contradict the
		equality of the half-relaxed limits with $u$.
	\end{proof}
	

	\begin{theorem}[Smooth-solution error bound]
		\label{thm:rate}
		Under Assumptions~\ref{ass:H1}--\ref{ass:H4}, let
		$u\in C_b^3(\mathcal N_T)$ be a classical solution of
		\eqref{eq:HJB-main} on a neighbourhood
		$\mathcal N_T\supset[0,T]\times\overline\cO$, where $C_b^3$ denotes
		bounded continuous derivatives of total order at most three.  Here
		``classical solution'' means that all three relations in
		\eqref{eq:HJB-main} hold pointwise; in particular,
		$\sup_{b\in B}q_b(t,x)=0$ on $[0,T)\times\partial\cO$, with $q_b$ defined
		below.  Set
		\[
		q_b(t,x):=\mathcal L(t,x,Du(t,x),u(t,x),b),
		\qquad (t,x,b)\in[0,T]\times\partial\cO\times B,
		\]
		and suppose that
		\begin{equation}
			\sup_{t\in[0,T],\,b\in B}
			\|q_b(t,\cdot)\|_{C^{1,1}(\partial\cO)}<\infty.
			\label{eq:residual-C11}
		\end{equation}
		Here $C^{1,1}(\partial\cO)$ is understood intrinsically on the compact
		hypersurface.  Fix a finite $C^2$ atlas
		$\{(U_\alpha,\kappa_\alpha)\}_\alpha$, with
		$\kappa_\alpha:U_\alpha\subset\partial\cO\to V_\alpha\subset\R^{N-1}$,
		and set
		\[
		\|v\|_{C^{1,1}(\partial\cO)}
		:=\max_\alpha
		\|v\circ\kappa_\alpha^{-1}\|_{C^{1,1}(V_\alpha)},
		\]
		where
		\[
		\|w\|_{C^{1,1}(V)}
		:=\|w\|_{L^\infty(V)}+\|Dw\|_{L^\infty(V)}
		+\operatorname{Lip}_V(Dw).
		\]
		Different finite atlases give equivalent norms.  In particular, if
		$v\in C^{1,1}(\partial\cO)$ attains a maximum at $x\in\partial\cO$, then,
		for $y\in\partial\cO$ sufficiently close to $x$,
		\begin{equation}
			v(y)\ge v(x)-C\|v\|_{C^{1,1}(\partial\cO)}|x-y|^2,
			\label{eq:C11-boundary-maximum}
		\end{equation}
		where $C$ depends only on the fixed atlas.
		Then the numerical solution satisfies
		\begin{equation}
			\max_{0\le n\le N_T}\max_{i\in\mathcal I_h}
			|U_i^n-u(t_n,x_i)|
			\le C\left(\dt^{1/2}+\frac{h^2}{\dt}\right).
			\label{eq:full-rate}
		\end{equation}
		The two terms on the right-hand side of~\eqref{eq:full-rate} are balanced
		by $\dt\asymp h^{4/3}$, which yields $O(h^{2/3})$; the commonly used
		choice $\dt\asymp h$ yields $O(h^{1/2})$.
	\end{theorem}

	\begin{proof}
		Set
		\[
		u_i^n:=u(t_n,x_i),\qquad
		u_h^n:=(u_j^n)_{j\in\mathcal I_h},\qquad
		e_i^n:=U_i^n-u_i^n,
		\]
		and
		\[
		E_n^\circ:=\max_{i\in\mathcal I_h^\circ}|e_i^n|,\qquad
		E_n^\partial:=\max_{i\in\mathcal I_h^\partial}|e_i^n|,\qquad
		E_n:=\max\{E_n^\circ,E_n^\partial\}.
		\]
		Here and below, $C$ is independent of $h,\dt,n,i,a,$ and $b$; it depends
		only on $T$, the coefficient, domain, and mesh bounds, the norm in
		\eqref{eq:residual-C11}, and $\|u\|_{C_b^3(\mathcal N_T)}$.
		
		\medskip
		\noindent\textbf{Interior nodes.}
		For $i\in\mathcal I_h^\circ$, let $\widehat{\cS}_{n,i}^{a,b}$ be the
		right-hand side of~\eqref{eq:fixed-operator} with the interpolation value
		$I_h[V](\widetilde X_p)$ replaced by the exact value
		$u(t_{n+1},\widetilde X_p)$, and set
		$\widehat S_{n,i}:=\inf_{a\in A,b\in B}\widehat{\cS}_{n,i}^{a,b}$.
		Adding and subtracting $S_{n,i}[u_h^{n+1}]$ and $\widehat S_{n,i}$ gives
		the exact decomposition
		\begin{equation}
			e_i^n=I_{1,i}^n+I_{2,i}^n+I_{3,i}^n,
			\label{eq:three-error-parts}
		\end{equation}
		where
		\begin{align*}
			I_{1,i}^n&:=\widehat S_{n,i}-u_i^n,\\
			I_{2,i}^n&:=S_{n,i}[U^{n+1}]-S_{n,i}[u_h^{n+1}],\\
			I_{3,i}^n&:=S_{n,i}[u_h^{n+1}]-\widehat S_{n,i}.
		\end{align*}
		
		The consistency calculation before the interpolation error is introduced
		gives
		\begin{align}
			I_{1,i}^n
			=-\sup_{a\in A,\,b\in B}\Biggl\{&
			\dt A_{n,i}^a
			+\frac1P\sum_{p=1}^P\widetilde d_{n,i}^p(a,b)
			\left(1-\frac12\widetilde k_{n,i}^p(a,b)
			\widetilde d_{n,i}^p(a,b)\right)
			q_b(t_n,x_b^p)\Biggr\}
			+O(\dt^{3/2}),
			\label{eq:I1-consistency}
		\end{align}
		where
		\[
		A_{n,i}^a:=
		\bigl[-\partial_tu+\mathcal H\bigr]
		(t_n,x_i,Du_i^n,D^2u_i^n,u_i^n,a),
		\]
		and a summand in~\eqref{eq:I1-consistency} is understood to be zero on a
		non-exiting branch.  Since $u$ satisfies the two classical equations,
		\[
		\sup_{a\in A}A_{n,i}^a=0,
		\qquad \sup_{b\in B}q_b(t,x)=0
		\quad (x\in\partial\cO),
		\]
		so every term inside the supremum in~\eqref{eq:I1-consistency} is
		nonpositive for sufficiently small $\dt$.
		
		We now bound that supremum from below.  Choose
		$a_*\in\arg\max_a A_{n,i}^a$.  If no branch for $a_*$ exits, the expression
		in braces vanishes for this $a_*$.  Otherwise, let $\bar x_i$ be a closest
		point of $x_i$ on $\partial\cO$ and choose
		$b_*\in\arg\max_b q_b(t_n,\bar x_i)$.  The step bound
		\eqref{eq:step-bound} first gives
		$\dist(x_i,\partial\cO)\le|X_{n,i}^{*,p_0}(a_*)-x_i|\le C\sqrt\dt$
		for any exiting branch $p_0$.  For every exiting branch generated by
		$(a_*,b_*)$, the definition of the oblique projection,
		Lemma~\ref{lem:oblique-distance}, and the triangle inequality give
		\[
		\begin{aligned}
			|x_{n,i}^{b_*,p}(a_*)-\bar x_i|
			&\le |x_{n,i}^{b_*,p}(a_*)-X_{n,i}^{*,p}(a_*)|
			+|X_{n,i}^{*,p}(a_*)-x_i|+|x_i-\bar x_i|\\
			&=\frac12\widetilde d_{n,i}^p(a_*,b_*)
			+|X_{n,i}^{*,p}(a_*)-x_i|
			+\dist(x_i,\partial\cO)
			\le C\sqrt\dt .
		\end{aligned}
		\]
		Moreover, $q_{b_*}(t_n,\cdot)\le0$ on $\partial\cO$ and attains its maximum
		zero at $\bar x_i$.  Hence its tangential gradient vanishes there, and the
		uniform $C^{1,1}$ bound~\eqref{eq:residual-C11} and
		\eqref{eq:C11-boundary-maximum} give
		\[
		-C|x_{n,i}^{b_*,p}(a_*)-\bar x_i|^2
		\le q_{b_*}(t_n,x_{n,i}^{b_*,p}(a_*))\le0.
		\]
		Since the supremum in~\eqref{eq:I1-consistency} is at least its value at
		the single pair $(a_*,b_*)$, the preceding estimate and
		$\widetilde d_{n,i}^p\le C\sqrt\dt$ prove
		\[
		-C\dt^{3/2}\le
		\sup_{a,b}\Biggl\{\dt A_{n,i}^a
		+\frac1P\sum_{p=1}^P\widetilde d_{n,i}^p
		\left(1-\frac12\widetilde k_{n,i}^p\widetilde d_{n,i}^p\right)
		q_b(t_n,x_b^p)\Biggr\}\le0,
		\]
		and therefore
		\begin{equation}
			|I_{1,i}^n|\le C\dt^{3/2}.
			\label{eq:I1-bound}
		\end{equation}
		
		Next, $|\inf F-\inf G|\le\sup|F-G|$, positivity of $I_h$, and
		$e^{-\widetilde k\widetilde d}/(1+r\dt)\le1$ give
		\begin{align}
			|I_{2,i}^n|
			&\le\sup_{a,b}
			\frac1{1+r(t_n,x_i,a)\dt}\frac1P\sum_{p=1}^P
			e^{-\widetilde k_{n,i}^p\widetilde d_{n,i}^p}
			I_h[|e^{n+1}|](\widetilde X_{n,i}^p)
			\le E_{n+1}.
			\label{eq:I2-bound}
		\end{align}
		The same inequality and the $\mathbb P_1$ interpolation estimate yield
		\begin{align}
			|I_{3,i}^n|
			&\le\sup_{a,b}
			\frac1{1+r(t_n,x_i,a)\dt}\frac1P\sum_{p=1}^P
			e^{-\widetilde k_{n,i}^p\widetilde d_{n,i}^p}
			\left|I_h[u_h^{n+1}](\widetilde X_{n,i}^p)
			-u(t_{n+1},\widetilde X_{n,i}^p)\right|
			\le Ch^2.
			\label{eq:I3-bound}
		\end{align}
		Thus the interior-node error satisfies
		\begin{equation}
			E_n^\circ\le E_{n+1}+C(\dt^{3/2}+h^2).
			\label{eq:smooth-interior-defect}
		\end{equation}
		
		\medskip
		\noindent\textbf{Boundary nodes.}
		Part~\textup{(II)} of Proposition~\ref{prop:nodewise-consistency} and the
		classical boundary condition give
		\begin{equation}
			\left|S_{n,i}^{\partial}[u_h^n]-u_i^n\right|\le Ch^2,
			\qquad i\in\mathcal I_h^\partial.
			\label{eq:boundary-exact-defect}
		\end{equation}
		Since $U_i^n=S_{n,i}^{\partial}[U^n]$ and
		$|\inf F-\inf G|\le\sup|F-G|$, it follows that
		\begin{equation}
			|e_i^n|\le
			\sup_{b\in B}
			\frac{\displaystyle\sum_{j\in\mathcal I_h^\partial}
				\omega_{ij}^b|e_j^n|
				+\displaystyle\sum_{j\in\mathcal I_h^\circ}
				\omega_{ij}^b|e_j^n|}
			{1+\ell_hk(t_n,x_i,b)}+Ch^2.
			\label{eq:boundary-error-direct}
		\end{equation}
		For small $h$, boundedness of $k$ and~\eqref{eq:boundary-row-margin} give
		\[
		\frac{\sum_{j\in\mathcal I_h^\circ}\omega_{ij}^b}
		{1+\ell_hk(t_n,x_i,b)}\ge\frac\theta2=: \theta_0
		\qquad(i\in\mathcal I_h^\partial,\ b\in B).
		\]
		If $E_n^\partial\le E_n^\circ$, no estimate is needed.  Otherwise, take a
		boundary node at which $E_n^\partial$ is attained.  Since the two normalized
		weight sums in~\eqref{eq:boundary-error-direct} have total at most one and
		the interior sum is at least $\theta_0$, we obtain
		\[
		E_n^\partial
		\le E_n^\partial-\theta_0(E_n^\partial-E_n^\circ)+Ch^2.
		\]
		Therefore
		\begin{equation}
			E_n^\partial\le E_n^\circ+\frac{C}{\theta_0}h^2.
			\label{eq:smooth-boundary-defect}
		\end{equation}
		Combining~\eqref{eq:smooth-interior-defect} and
		\eqref{eq:smooth-boundary-defect} only at this point yields
		\[
		E_n\le E_{n+1}+C(\dt^{3/2}+h^2).
		\]
		Since $E_{N_T}=0$ and $N_T\dt=T$, backward iteration gives
		\[
		E_n\le C(N_T-n)(\dt^{3/2}+h^2)
		\le CT\left(\dt^{1/2}+\frac{h^2}{\dt}\right),
		\]
		which is~\eqref{eq:full-rate}.
	\end{proof}

\section{Numerical experiments}
\label{sec:numerics}

We report four two-dimensional tests for the split scheme~\eqref{eq:scheme}.
The MATLAB implementation precomputes the positive $\mathbb P_1$ rows and all
time-independent branch geometry, retains two time layers, and factorizes each
same-level boundary system once.  The tables contain errors and successive
observed orders; timings and algebraic diagnostics are retained in the
accompanying data files.

\paragraph{Common discretization.}
Experiments~1--3 use the same five nested fitted triangulations of the unit
disk.  Their nominal dyadic parameters are
$h_{\rm nom}=2^{-3},\ldots,2^{-7}$, with respectively
$377$, $1441$, $5633$, $22273$, and $88577$ nodes.  The realized maximum edge
length is $h_{\rm real}=1.243232\,h_{\rm nom}$ on every level; hence the
successive orders are unchanged whether $h_{\rm nom}$ or $h_{\rm real}$ is
used.  Experiment~1 uses $256$ equally spaced unit-circle controls, excluding
the zero control, in both compared implementations.  The main refinement in
Experiment~2 uses the zero control together with $256$ circle directions; its
control sensitivity uses $K=32,64,128,256$ circle directions, always with the
zero control retained.  Experiments~3--4 use the zero control together with
$64$ circle directions.  Both weak Euler branches have probability $1/2$.
Smooth-disk probes are evaluated by the same positive affine-chord
$\mathbb P_1$ map.  Experiment~4 uses its own fitted mixed-boundary meshes.
We set $\ell_h=h_{\rm nom}$.

The angular control error is kept separate from the space--time refinement.
For
\[
 A_K:=\{(\cos(2\pi j/K),\sin(2\pi j/K)):0\le j<K\},
 \qquad H_K(p):=\max_{a\in A_K}a\cdot p,
\]
the nearest angular direction gives
\begin{equation}
 0\le |p|-H_K(p)
 \le \eta_K|p|,
 \qquad \eta_K:=1-\cos(\pi/K).
 \label{eq:angular-control-error}
\end{equation}
For the even values of $K$ used below, adding the zero control does not change
$H_K$, because $A_K$ contains antipodal pairs and hence $H_K(p)\ge0$.
Here $\eta_{64}=1.20454\times10^{-3}$ and
$\eta_{256}=7.52982\times10^{-5}$.  The smooth manufactured solution used in
Experiments~1 and~2 satisfies $\|Du\|_\infty\le3/2$, so the Hamiltonian
residual is bounded by
$1.80682\times10^{-3}$ for $K=64$ and by
$1.12947\times10^{-4}$ for $K=256$.
For the present method, replacing $A$ by $A_K$ adds at most
$\Delta t\,\eta_K\|Du\|_\infty$ to the analytic one-step defect; the recursion
in Theorem~\ref{thm:rate} therefore adds $C\eta_K$ to the full-grid bound.
Thus the tables report the combined space--time and finite-control errors over
the displayed range.  With fixed $K$ an angular-error plateau may eventually
occur; joint asymptotic convergence to the continuous-control problem
requires $K=K(h)\to\infty$.

Every time grid ends at $T=1$: for a prescribed nominal step,
\begin{equation}
 N_T=\operatorname{round}(T/\Delta t_{\rm nom}),
 \qquad \Delta t=T/N_T.
 \label{eq:numerical-time-alignment}
\end{equation}
At $t=0$ we use
\[
 E_\infty=\max_i|U_i-u(0,x_i)|
\]
and, if $\widehat K$ is a curved auxiliary cell with geometric barycentre
$x_{\widehat K}$,
\begin{equation}
 E_1=\sum_{\widehat K}|\widehat K|
 \left|I_h[U](x_{\widehat K})-u(0,x_{\widehat K})\right|.
 \label{eq:barycentric-E1}
\end{equation}
For two consecutive dyadic levels,
\[
 p=\frac{\log(E_{\rm old}/E_{\rm new})}
          {\log(h_{{\rm nom,old}}/h_{{\rm nom,new}})}.
\]
The first displayed level is assigned ``--''.

\subsection{Experiment 1: calibrated \texorpdfstring{$k=0$}{k=0}
comparison on the unit disk}
\label{sec:disk-comparison}

Let
\[
 u(t,x_1,x_2)=(t+0.5)\sin x_1\sin x_2,\qquad
 \sigma(x)=\sqrt2
 \begin{pmatrix}\sin(x_1+x_2)\\ \cos(x_1+x_2)\end{pmatrix}.
\]
With $\mu(a)=-a$, $|a|\le1$, and $r=k=0$, set
\begin{equation}
 f=-u_t-\frac12\sigma^\top D^2u\,\sigma+|Du|.
 \label{eq:experiment2-source}
\end{equation}
On the boundary $g=\gamma\cdot Du$, first with $\gamma=n$ and then with
$\gamma=R_{-\pi/6}n$, where $R_\alpha$ denotes counterclockwise rotation
through $\alpha$.  The reference benchmark is compared after the
corresponding exact time reversal.

For comparison we independently implement the all-node formulas
(25)--(30) of~\cite{CalzolaEtAl2023}; this implementation has no separate
same-level boundary solve.  The present and reference implementations use the
same meshes, time grids, controls, boundary directions, interpolation map, and
error definitions.  We fix $\Delta t=h_{\rm nom}$.  The additional inward
offset $\bar c$ belongs only to the reference formula and is selected using
only errors of that implementation.  On the calibration levels
$h_{\rm nom}=2^{-3},2^{-4},2^{-5}$, we first test
$\bar c=0.050,0.075,\ldots,1.500$.  The score of each candidate is the minimum
of the eight successive calibration orders obtained from two refinements, two
boundary directions, and the two error norms.  This scan selects
$\bar c=0.675$, with minimum and mean scores $0.943$ and $0.979$.
A prespecified refinement over $\bar c=0.650,0.655,\ldots,0.700$ then selects
\begin{equation}
 \bar c=0.690,
 \qquad \min p_{\rm cal}=0.945,
 \qquad \operatorname{mean}p_{\rm cal}=0.979.
 \label{eq:selected-cbar}
\end{equation}
The selected value is frozen before the two finer levels

\begin{thebibliography}{99}
\bibitem[Barles1993]{Barles1993} G.~Barles, Fully nonlinear Neumann type boundary conditions for second-order elliptic and parabolic equations, \emph{Journal of Differential Equations} \textbf{106} (1993), no.~1, 90--106.
\bibitem[Bourgoing2008]{Bourgoing2008} M.~Bourgoing, Viscosity solutions of fully nonlinear second order parabolic equations with $L^1$ dependence in time and Neumann boundary conditions, \emph{Discrete and Continuous Dynamical Systems} \textbf{21} (2008), no.~3, 763--800.
\bibitem[CalzolaEtAl2023]{CalzolaEtAl2023} E.~Calzola, E.~Carlini, X.~Dupuis, and F.~J.~Silva, A semi-Lagrangian scheme for Hamilton--Jacobi--Bellman equations with oblique derivatives boundary conditions, \emph{Numerische Mathematik} \textbf{153} (2023), 49--84.
\bibitem[CrandallIshiiLions1992]{CrandallIshiiLions1992} M.~G.~Crandall, H.~Ishii, and P.-L.~Lions, User's guide to viscosity solutions of second order partial differential equations, \emph{Bulletin of the American Mathematical Society} \textbf{27} (1992), 1--67.
\bibitem[GilbargTrudinger2001]{GilbargTrudinger2001} D.~Gilbarg and N.~S.~Trudinger, \emph{Elliptic Partial Differential Equations of Second Order}, Springer-Verlag, Berlin, 2001.
\bibitem[Gobet2001]{Gobet2001} E.~Gobet, Efficient schemes for the weak approximation of reflected diffusions, \emph{Monte Carlo Methods and Applications} \textbf{7} (2001), no.~1--2, 193--202.
\end{thebibliography}
\end{document}
